\documentclass[11pt]{article}

\usepackage[margin=1in]{geometry}
\usepackage{amsmath,amssymb,amsthm,mathtools,bm}
\usepackage{booktabs}
\usepackage{enumitem}
\usepackage{microtype}
\usepackage{xcolor}
\usepackage[hidelinks]{hyperref}
\hypersetup{
  pdftitle={Directed Reachability from Cut Queries in
            O-tilde(n to the three-halves) Queries},
  pdfauthor={Anonymous Authors},
  pdfsubject={Directed graph algorithms, cut-query complexity, and
              implicit interior-point methods},
  pdfkeywords={directed reachability, cut queries, spectral sparsification,
               interior-point methods, Lewis weights}
}

\allowdisplaybreaks
\setlist{nosep}

\newtheorem{theorem}{Theorem}[section]
\newtheorem{lemma}[theorem]{Lemma}
\newtheorem{corollary}[theorem]{Corollary}
\newtheorem{proposition}[theorem]{Proposition}
\newtheorem{claim}[theorem]{Claim}
\theoremstyle{definition}
\newtheorem{definition}[theorem]{Definition}
\newtheorem{remark}[theorem]{Remark}

\newcommand{\R}{\mathbb{R}}
\newcommand{\eps}{\varepsilon}
\newcommand{\poly}{\operatorname{poly}}
\newcommand{\polylog}{\operatorname{polylog}}
\newcommand{\rank}{\operatorname{rank}}
\newcommand{\diag}{\operatorname{diag}}
\newcommand{\supp}{\operatorname{supp}}
\newcommand{\OPT}{\mathrm{OPT}}
\newcommand{\id}{\mathbf{1}}
\newcommand{\ones}{\mathbf{1}}
\newcommand{\E}{\mathbb{E}}
\newcommand{\Prb}{\mathbb{P}}
\newcommand{\cH}{c_H}
\newcommand{\Lap}{\mathcal{L}}
\newcommand{\norm}[1]{\left\lVert #1\right\rVert}
\newcommand{\abs}[1]{\left\lvert #1\right\rvert}
\newcommand{\set}[1]{\left\{#1\right\}}
\newcommand{\ip}[2]{\left\langle #1,#2\right\rangle}
\newcommand{\wtO}{\widetilde O}
\newcommand{\wtOmega}{\widetilde\Omega}
\newcommand{\defeq}{\mathrel{:=}}

\title{Directed Reachability from Cut Queries in
       $\widetilde O(n^{3/2})$ Queries}

\date{}

\begin{document}
\maketitle

\begin{abstract}
We study directed $s$--$t$ reachability when an unknown digraph
$G=(V,E)$ is accessible only through an outgoing-cut oracle: a query
$S\subseteq V$ returns $F(S)=|E(S,V\setminus S)|$.  We give a randomized
adaptive algorithm using $n^{3/2}\log^{O(1)}n$ queries and succeeding with
probability $1-n^{-\Omega(1)}$; computation between queries is unrestricted.
The conceptual core is a rank-$n$ fractional-orientation linear program.  An
inverse-polynomially accurate solution induces an implicit threshold digraph
having exactly the same zero out-cuts, and hence the same reachability
relation, as $G$.  Although the LP has up to $\Theta(n^2)$ hidden rows, we
implement a rank-sensitive Lee--Sidford box-LP method without materializing
them.  Our main oracle primitive is a spread-independent weighted spectral
sparsifier for pair-evaluable conductances, constructed from cut queries by
filtered weighted spanners, bundle leverage bounds, and whole-pair sampling.
Duplicate hidden copies are kept block-symmetric.  Approximate normal solves
are handled by the standard stable Lee--Sidford interface, and a public star
right inverse keeps every primal update exactly feasible.  The result
breaks the quadratic query barrier for directed reachability posed in recent
work on the cut-query model.
\end{abstract}

\textbf{Note:} The algorithm and main proof were discovered by an agentic harness developed by George Z. Li and Janardhan Kulkarni that is built on OpenAI’s GPT-5.6 Sol and Gemini 3.1 Pro. 

\section{Introduction}

Let $G=(V,E)$ be an unknown directed graph on $n$ labeled vertices.  For a
set $S\subseteq V$, define its outgoing cut value by
\[
    F(S) \defeq |E(S,V\setminus S)|.
\]
The algorithm may adaptively query $F$ on arbitrary sets, but receives no
other information about $E$.  Given $s,t\in V$, the goal is to decide whether
$G$ contains a directed path from $s$ to $t$ while minimizing the number of
queries.  We do not charge computation between queries.

Directed reachability is a particularly stark test of the power of this
oracle.  The direct reconstruction algorithm uses $O(n^2)$ pair tests, and
the outgoing-cut function does not determine the orientation of each arc:
reversing a directed cycle leaves every cut value unchanged.  Nevertheless,
$F$ directly identifies all zero out-cuts, and those cuts characterize
reachability.  The question of obtaining any
$O(n^{2-\gamma})$ bound, for a constant $\gamma>0$, was highlighted as open
in recent work on cut queries~\cite{JiangNanongkaiSawettamalya2026} and as an
explicit open problem by Nanongkai~\cite{NanongkaiOpen}.

Our main result gives an exponent $3/2$.

\begin{theorem}[Main theorem]\label{thm:main}
Let $G=(V,E)$ be a loopless simple directed graph on $n$ vertices (so there
is at most one copy of each ordered arc), and let $s,t\in V$.  There is a
randomized adaptive algorithm which, given oracle access to
$F(S)=|E(S,V\setminus S)|$, decides whether $t$ is reachable from $s$,
succeeds with probability $1-n^{-\Omega(1)}$, and makes
\[
        n^{3/2}\log^{O(1)}n
\]
oracle queries.  More generally, if $G$ is a loopless directed multigraph
and the total number of arc copies in both directions on every unordered
pair is at most a constant $C$, the same algorithm uses
$O(Cn^{3/2}\log^{O(1)}n)$ queries.  (The simple-graph case has $C\le2$.)
Computation between queries is unrestricted, and the displayed query bound
is a deterministic cap for every random tape.
\end{theorem}

Self-loops never affect a cut or reachability between distinct vertices and
may be discarded.  The case $s=t$ is immediate.

\subsection{Overview of the proof}

The proof has four layers.

\paragraph{Cut algebra.}
For each unordered pair $\{u,v\}$, let
\[
 c_{uv}=\id[(u,v)\in E]+\id[(v,u)\in E]\in\{0,1,2\},
\]
and let $H$ be the undirected multigraph with $c_{uv}$ copies of
$\{u,v\}$.  The imbalance
\[
 b_v=F(\{v\})-F(V\setminus\{v\})
\]
is the outdegree minus the indegree of $v$.  Direct cancellation gives
\[
        2F(S)=\cH(\delta S)+b(S).
\]
Thus the oracle exposes both the cut function of $H$ and the imbalance
vector $b$, even though it does not expose the individual directions.

\paragraph{A fractional orientation certificate.}
Reference-orient every pair $u<v$ from $u$ to $v$ and associate a variable
$z_{uv}\in[-1,1]$.  It represents the fractional orientation
\[
 p_{u\to v}=\frac{1+z_{uv}}2,
 \qquad p_{v\to u}=\frac{1-z_{uv}}2.
\]
We construct $z$ whose divergence is extremely close to $(1-\eta)b$, where
$\eta=2^{-\lceil12\log_2n\rceil}\le n^{-12}$.  If $P_z(S)$ is the total fractional mass
leaving $S$, then
\[
  2P_z(S)=2(1-\eta)F(S)+\eta\cH(\delta S)+\rho(S),
  \qquad \norm{\rho}_1\le n^{-20}.
\]
It follows that $P_z(S)<n^{-4}$ when $F(S)=0$, whereas
$P_z(S)>1/2$ when $F(S)>0$.  Thresholding individual fractional arcs at
$n^{-4}$ therefore preserves every zero out-cut and hence every reachability
relation.

\paragraph{A hidden rank-$n$ box LP.}
We obtain $z$ from a Phase-I box LP with one conceptual row per copy in $H$
and two public star rows per vertex.  It may have $\Theta(n^2)$ rows but has
rank $n$.  Its optimum is zero: the unknown original orientation itself gives
a block-symmetric witness.  We apply the rank-sensitive box-LP framework of
Lee and Sidford~\cite{LeeSidford2020}, which uses
$\sqrt n\log^{O(1)}n$ weighted-system-equivalent calls.

\paragraph{Implicit implementation.}
Each call is implemented in $n\log^{O(1)}n$ cut queries.  A new weighted
spectral primitive constructs a sparsifier of
\[
  \sum_{u<v}c_{uv}D_{uv}(e_u-e_v)(e_u-e_v)^\top
\]
whenever $D_{uv}$ is computable from the endpoints.  It yields both a
constant-condition preconditioner and simultaneous leverage estimates for
all hidden rows.  Bit-class decomposition implements all coordinate
aggregates.  A KKT formula implements the mixed-norm support oracle.  Finally,
an explicit star right inverse repairs each approximate direction exactly, so
feasibility never drifts.  Multiplying the call count by the
per-call query cost gives
\[
  \sqrt n\log^{O(1)}n\cdot n\log^{O(1)}n
       =n^{3/2}\log^{O(1)}n.
\]

\paragraph{How to read the proof.}
The graph-theoretic correctness spine is
Sections~\ref{sec:prelim}, \ref{sec:orientation}, and
\ref{sec:completion}: cut algebra exposes an undirected support, the
fractional certificate preserves zero out-cuts, and the final section turns
an approximate LP solution into a reachability answer.  On a first reading,
Theorem~\ref{thm:implicit-lp} may be treated as a black box while following
that spine.  Section~\ref{sec:spectral} constructs the spectral oracle used
to realize the black box; Section~\ref{sec:implicit-lp} implements its
path-following primitives; and Section~\ref{sec:numerics} checks the short
interface between those cut-query primitives and the stable Lee--Sidford
solver.  Thus a
reader interested first in the reduction need not understand interior-point
methods before reading the correctness argument.

\paragraph{Core notation.}
The symbols used across several layers are collected here for reference.
\begin{center}
\small
\begin{tabular}{@{}lp{0.77\linewidth}@{}}
\toprule
Symbol & Meaning\\
\midrule
$F(S)$ & number of input arc copies directed from $S$ to $V\setminus S$\\
$c_{uv},H$ & total multiplicity on $\{u,v\}$, and the resulting undirected multigraph\\
$b$ & input outdegree-minus-indegree vector; $d=(1-\eta)b$ is the Phase-I right-hand side\\
$a_e$ & reference-oriented incidence vector $e_u-e_v$ for $e=\{u,v\}$, $u<v$\\
$z_e,p_{u\to v}$ & common hidden-pair variable and its induced fractional directed mass\\
$J,\tau$ & threshold digraph and threshold $\tau=n^{-4}$\\
$A,x,M$ & full Phase-I row matrix, its row-variable vector, and its number of rows\\
\bottomrule
\end{tabular}
\end{center}
Section~\ref{sec:implicit-lp} gives a second, solver-state table before
introducing the IPM notation.

\subsection{Our contributions}

Beyond the main bound, the proof provides the following reusable ingredients.
\begin{enumerate}[label=(\roman*)]
\item A reduction from outgoing-cut reachability to a rank-$n$
      fractional-orientation LP whose approximate solutions preserve the
      complete zero-outcut family.
\item A spread-independent spectral sparsifier for unknown weighted
      multigraphs with pair-evaluable conductances, using
      $O(\eps^{-2}n\log^8(n/(\eps\delta)))$ cut queries.
\item A block-symmetric cut-query implementation of a rank-sensitive box-LP
      method, including simultaneous leverage estimation without assigning
      independent random signs to indistinguishable parallel copies.
\item A short interface from cut-query primitives to the standard stable
      Lee--Sidford execution: residual-certified Richardson iteration gives
      inverse-polynomial solve accuracy, and a public star right inverse
      preserves affine feasibility exactly.
\end{enumerate}

\subsection{Related work}

Rubinstein, Schramm, and Weinberg~\cite{RubinsteinSchrammWeinberg2018}
introduced the modern cut-query study of exact graph cuts and obtained
$\widetilde O(n^{5/3})$ queries for undirected minimum $s$--$t$ cut.
Subsequent work obtained linear-query undirected edge connectivity via star
contraction~\cite{ApersEtAl2022}, deterministic
$\widetilde O(n^{5/3})$ undirected flow and connectivity bounds
~\cite{AnandSaranurakWang2025}, and an
$\widetilde O(n^{8/5})$ randomized bound for undirected minimum $s$--$t$
cut~\cite{JiangNanongkaiSawettamalya2026}.  The last work explicitly notes
that no subquadratic upper bound was known even for directed reachability.
Our oracle is asymmetric: it returns arcs leaving $S$, rather than undirected
edges crossing $S$.

For linear programming we use the rank-sensitive box-LP framework of Lee and
Sidford~\cite{LeeSidford2020}.  Its approximate Lewis-weight interface is
part of that framework; the stability analysis for inverse-polynomially
accurate linear solves is given in their earlier primal path-following
work~\cite[Section~8]{LeeSidfordPathII}.  We use those results as a solver
interface rather than reproving path-following stability.

Our spectral routine builds on weighted spanners of Baswana and
Sen~\cite{BaswanaSen2007}, the bundle-spanner sampling paradigm of
Koutis~\cite{Koutis2014}, and matrix Bernstein inequalities in the form
developed by Tropp~\cite{Tropp2012}.  The new points here are the filtered
cut-query simulation, frozen phase snapshots, multiplicity-block sampling,
and an interface for endpoint-computable weights of arbitrary spread.

\section{Preliminaries and cut-query algebra}\label{sec:prelim}

For an undirected multigraph $H$, let $\cH(\delta S)$ denote the number of
copies crossing $(S,V\setminus S)$, and let $\cH(A,B)$ denote the number of
copies with one endpoint in each of two disjoint sets $A,B$.  We use
$L\preceq M$ for the Loewner order on symmetric matrices.  For a positive
semidefinite matrix $L$, $L^\dagger$ is its Moore--Penrose pseudoinverse.
The effective resistance of $e=uv$ in a weighted graph with Laplacian $L$ is
$(e_u-e_v)^\top L^\dagger(e_u-e_v)$.

\begin{definition}[Pair-evaluable scalar]
A scalar $q_{uv}$ is \emph{pair-evaluable} if, from the public transcript and
the labels $u,v$, the algorithm can compute $q_{uv}$ without querying whether
$c_{uv}>0$.  Offline computation and the size of the representing arithmetic
circuit are not charged.
\end{definition}

This definition allows dependence on public hashes, current endpoint states,
known cluster labels, and all previous answers.

\paragraph{Query-model convention (optional on a first reading).}
As is standard for query complexity, exact real arithmetic and comparisons
between queries are free.  A primal coordinate or a path weight may therefore
be represented by a public arithmetic circuit rather than by a fixed-precision
word.  Whenever such a value enters a cut-query reduction, certified
evaluation freezes a dyadic approximation of the accuracy requested by the
solver.  Lemma~\ref{lem:rounded-aggregate} bounds the resulting error, and
Section~\ref{sec:numerics} checks that the stable Lee--Sidford interface
absorbs it.

Filtered search may evaluate a row formula on an absent candidate pair.  We
make every such circuit total by clipping virtual slacks away from zero; this
choice is irrelevant to an aggregate because an absent pair is multiplied by
$c_{uv}=0$.  Before using a batch, one filtered search rejects if an
\emph{actual} row is outside the domain required by the original formula.
The Lee--Sidford interiority and conditioning bounds imply that this never
happens on a successful execution.  Fixed caps make an unsuccessful random
execution terminate within the advertised query bound.

\begin{lemma}[Cut algebra]\label{lem:cut-algebra}
Let
\[
 c_{uv}=\id[(u,v)\in E]+\id[(v,u)\in E]
 \quad\text{and}\quad
 b_v=F(\{v\})-F(V\setminus\{v\}).
\]
Let $H$ be the undirected multigraph with multiplicities $c_{uv}$.  Then for
every $S\subseteq V$,
\begin{equation}\label{eq:cut-algebra}
        2F(S)=\cH(\delta S)+b(S).
\end{equation}
Moreover, for disjoint $A,B\subseteq V$,
\begin{equation}\label{eq:rectangle}
        \cH(A,B)=F(A)+F(B)-F(A\cup B).
\end{equation}
\end{lemma}

\begin{proof}
For every arc internal to $S$ or to $V\setminus S$, both sides of
\eqref{eq:cut-algebra} receive zero.  An arc from $S$ to $V\setminus S$
contributes two to the right-hand side: one through $\cH(\delta S)$ and one
through $b(S)$.  An arc in the opposite direction contributes one and minus
one, hence zero.  This proves \eqref{eq:cut-algebra}.  In
$F(A)+F(B)-F(A\cup B)$, all arcs from $A\cup B$ to its complement cancel;
the surviving arcs are precisely those in either direction between $A$ and
$B$, proving \eqref{eq:rectangle}.
\end{proof}



\begin{corollary}[Filtered-neighbor oracle]\label{cor:filtered-neighbor}
Fix $u\notin T$.  Three outgoing-cut queries determine
$\sum_{v\in T}c_{uv}$ for every publicly computable candidate set $T$.
Consequently, a qualifying neighbor can be found by binary search in
$O(\log n)$ such tests.
\end{corollary}

The candidate set may depend on the entire public transcript but not on an
unqueried adjacency bit.  This distinction is used repeatedly below.

\begin{lemma}[Bit-class aggregation]\label{lem:aggregate}
Let $q_{uv}$ be a signed $B$-bit pair-evaluable scalar.  All values
\[
        Q_u=\sum_{v\ne u}c_{uv}q_{uv},\qquad u\in V,
\]
can be computed using $O(nB)$ outgoing-cut queries.  The same bound applies
to signed incidence products, predicate-restricted sums, and weighted
Laplacian products whose coordinate data are $B$-bit pair circuits.
\end{lemma}

\begin{proof}
Write the positive and negative parts of $q_{uv}$ in binary.  For a fixed
$u$, sign, and bit position $j$, form the public set of vertices whose
corresponding bit is one and apply Corollary~\ref{cor:filtered-neighbor}.
Recombining the answers gives $Q_u$ exactly for the stored fixed-point
values.  Incidence signs and public predicates are incorporated when the
candidate sets are formed.
\end{proof}

\paragraph{Example: reconstructing one weighted sum.}
Fix $u$ and suppose three candidate endpoints $v_1,v_2,v_3$ have nonnegative
three-bit values
\[
 q_{uv_1}=5=(101)_2,\qquad
 q_{uv_2}=2=(010)_2,\qquad
 q_{uv_3}=3=(011)_2.
\]
The public endpoint sets for the $4$-, $2$-, and $1$-bits are respectively
\[
 T_2=\{v_1\},\qquad T_1=\{v_2,v_3\},\qquad
 T_0=\{v_1,v_3\}.
\]
Using the rectangle identity to obtain
$\cH(\{u\},T_j)=\sum_{v\in T_j}c_{uv}$, the algorithm reconstructs
\[
 \sum_{i=1}^3 c_{uv_i}q_{uv_i}
 =4\cH(\{u\},T_2)+2\cH(\{u\},T_1)+\cH(\{u\},T_0).
\]
Signed values are handled by performing the same calculation separately for
their positive and negative parts.  Notice that the sets $T_j$ depend only
on the pair-evaluable values; they do not reveal which candidate pairs
actually occur.

\begin{lemma}[Rounded circuit aggregation]\label{lem:rounded-aggregate}
Let $q_{uv}$ be a pair-evaluable circuit and suppose a public envelope gives
$|q_{uv}|\le\Lambda$.  If certified evaluation produces a signed dyadic
$\widehat q_{uv}$ with
\[
        |\widehat q_{uv}-q_{uv}|\le\beta ,
\]
then all sums
\[
 \widehat Q_u=\sum_{v\ne u}c_{uv}\widehat q_{uv}
\]
are computed exactly with
$O(n\log(\Lambda/\beta))$ outgoing-cut queries and satisfy
\[
 |\widehat Q_u-Q_u|
 \le \beta\sum_{v\ne u}c_{uv}.
\]
If $D$ and $\widehat D$ are diagonal row weights with
$|D_i-\widehat D_i|\le\beta$, every row of $A$ has squared norm at most two,
and $A$ has $M$ rows, then
\begin{equation}\label{eq:rounded-normal-error}
 \norm{A^\top(\widehat D-D)A}_2\le2M\beta .
\end{equation}
The same conclusions hold for predicate-restricted sums and right-hand sides.
\end{lemma}

\begin{proof}
Apply Lemma~\ref{lem:aggregate} to the certified dyadic values.  The scalar
error bound follows by the triangle inequality.  For the matrix statement,
\[
 \norm{\sum_i(\widehat D_i-D_i)a_ia_i^\top}_2
 \le\sum_i\beta\norm{a_i}_2^2\le2M\beta .
\]
\end{proof}

\begin{lemma}[Hidden maxima and predicates]\label{lem:hidden-max}
Let $q_{uv}$ be a public pair-evaluable dyadic value and let $P(u,v)$ be a
public predicate.  The maximum of $q_{uv}$ over actual support pairs satisfying
$P$, an actual maximizing pair when one exists, and emptiness of the
qualifying support can be found using $O(n\log n)$ outgoing-cut queries.
\end{lemma}

\begin{proof}
For each fixed $u$, order all candidate endpoints $v$ by
$(q_{uv},v)$.  Binary search this public order, using
Corollary~\ref{cor:filtered-neighbor} on suffixes intersected with $P$, to find
the largest actual qualifying neighbor of $u$.  Take the best of the at most
$n$ returned candidates offline.  The same search with a constant value
implements emptiness and threshold predicates.
\end{proof}

\section{A fractional-orientation certificate}\label{sec:orientation}

Reference-orient each unordered pair $e=\{u,v\}$ from the smaller endpoint
$u$ to the larger endpoint $v$, and write $a_e=e_u-e_v\in\R^n$.  Each of the
$c_e$ conceptual copies receives a variable $z_{e,j}$.  The implementation
will keep all copies equal, but it is useful to retain the conceptual copies
when defining the LP.

\paragraph{What the Phase-I construction is doing.}
The reference sign fixes the meaning of every hidden coordinate: a positive
$z_{uv}$ contributes positive divergence at the smaller endpoint $u$ and
negative divergence at $v$.  There are $c_{uv}$ conceptual coordinates
because an unordered pair may support several input arc copies.  They are
kept equal by the implementation, but retaining the copies makes their total
incidence contribution exactly $c_{uv}a_{uv}z_{uv}$.  The extra grounded
root has no retained coordinate, so a star edge from $v$ to the root has row
$e_v$ (and its reverse has row $-e_v$).  The difference
$y_v^+-y_v^-$ can therefore absorb any remaining divergence at $v$.
Minimizing $y_v^++y_v^-$ drives the $\ell_1$ size of this public correction
to zero.  Finally, the factors $1-\eta$ and $\eta/2$ keep the optimum away
from the hidden box boundary, which supplies the strict interior required by
the barrier method.

\paragraph{A three-vertex example.}
Suppose $\{1,2\}$ supports only the arc $1\to2$, while $\{2,3\}$ supports one
arc in each direction.  Then
\[
 c_{12}=1,\quad \Delta_{12}=1,\qquad
 c_{23}=2,\quad \Delta_{23}=0,
\]
and the reference rows are $a_{12}=e_1-e_2$ and
$a_{23}=e_2-e_3$.  The block-symmetric zero-objective witness used below
sets
\[
 z_{12,1}=1-\eta,\qquad z_{23,1}=z_{23,2}=0.
\]
Thus the single directed pair supplies its shrunken original orientation,
whereas the antiparallel pair has zero net divergence.  At the known
starting point all hidden variables are instead zero, and the two explicit
star variables at each vertex supply the entire target divergence $d$.

Set
\[
        \eta=2^{-\lceil12\log_2n\rceil}\le n^{-12},
        \qquad d=(1-\eta)b,
        \qquad K_0=1+\max_v |d_v|.
\]
Add a new root and ground its coordinate.  For every original vertex $v$,
add two explicit rows $e_v$ and $-e_v$, with variables $y_v^+$ and $y_v^-$.
Consider the closed-box LP
\begin{equation}\label{eq:phase-one}
\begin{aligned}
 \min\quad & \sum_{v\in V}(y_v^++y_v^-) \\
 \text{s.t.}\quad &
   \sum_{e}\sum_{j=1}^{c_e}a_e z_{e,j}
      +\sum_{v\in V}e_v(y_v^+-y_v^-)=d,\\
 & -1+\eta/2\le z_{e,j}\le1-\eta/2,\\
 & 0\le y_v^+,y_v^-\le2K_0.
\end{aligned}
\end{equation}
The barrier algorithm operates in the strict interior of these boxes.
If $B\in\R^{m\times n}$ contains one row $a_e^\top$ per conceptual hidden
copy, then the full row matrix and primal vector are
\[
 M=m+2n,\qquad
 A=\begin{bmatrix}B\\I_n\\-I_n\end{bmatrix}
    \in\R^{M\times n},
 \qquad x=(z,y^+,y^-)\in\R^M,
\]
and the equality in \eqref{eq:phase-one} is $A^\top x=d$.
The conceptual hidden row count is also known without reconstructing the
support:
\[
        m=|E|=\sum_{v\in V}F(\{v\}),
\]
because every directed arc is counted once at its tail.

\begin{lemma}[Phase-I structure]\label{lem:phase-one}
The constraint matrix of \eqref{eq:phase-one} has rank $n$.  The point
\[
 z_{e,j}=0,\qquad
 y_v^+=K_0+d_v/2,\qquad y_v^-=K_0-d_v/2
\]
is explicitly known, exactly feasible, and strictly interior.  Moreover,
the optimum of \eqref{eq:phase-one} is zero, and there is an optimal hidden
assignment constant across all copies of an unordered pair.
\end{lemma}

\begin{proof}
The $e_v$ rows form an identity matrix, so the rank is $n$.  The displayed
point has signed star flow $d$ and zero hidden flow.  Since $K_0>|d_v|$, both
star variables lie strictly between $0$ and $2K_0$.

For $e=\{u,v\}$ with $u<v$, define
\[
 \Delta_e=\id[(u,v)\in E]-\id[(v,u)\in E].
\]
For every copy set
\[
        z_{e,j}=\frac{(1-\eta)\Delta_e}{c_e},
        \qquad y_v^+=y_v^-=0.
\]
Its divergence is $(1-\eta)b=d$.  Because $|\Delta_e|\le c_e$, each hidden
coordinate lies in $[-1+\eta,1-\eta]$, strictly inside the narrower hidden
box.  If both antiparallel arcs exist, then $c_e=2$ and $\Delta_e=0$, so the
common-copy assignment remains valid.  The objective is nonnegative and the
displayed witness has value zero.
\end{proof}

Suppose now that a common pair value $z_e$ satisfies $|z_e|\le1$.  For
$e=\{u,v\}$, $u<v$, define
\begin{equation}\label{eq:fractional-orientation}
 p_{u\to v}=\frac{1+z_e}{2},
 \qquad
 p_{v\to u}=\frac{1-z_e}{2}.
\end{equation}
For an arbitrary ordered pair, this means
\[
 p_{u\to v}=\begin{cases}
 (1+z_{uv})/2,&u<v,\\
 (1-z_{vu})/2,&v<u.
 \end{cases}
\]
Let $P_z(S)$ be the sum, over all conceptual copies crossing the cut, of the
fractional mass directed from $S$ to $V\setminus S$.

\begin{lemma}[Zero/positive outcut gap]\label{lem:gap}
Assume
\begin{equation}\label{eq:certificate}
 |z_e|\le1\quad\text{for every actual pair},
 \qquad
 \norm{B^\top z-(1-\eta)b}_1\le n^{-20},
\end{equation}
where $B^\top z=\sum_e c_ea_ez_e$.  Put
$\rho=B^\top z-(1-\eta)b$.  Then, for every $S\subseteq V$,
\begin{equation}\label{eq:fractional-cut}
 2P_z(S)=2(1-\eta)F(S)+\eta\cH(\delta S)+\rho(S).
\end{equation}
Consequently, for all sufficiently large $n$,
\[
 F(S)=0 \Longrightarrow P_z(S)<n^{-4},
 \qquad
 F(S)>0 \Longrightarrow P_z(S)>1/2.
\]
\end{lemma}

\begin{proof}
For any fractional orientation, twice the outgoing mass is the undirected
crossing multiplicity plus the fractional divergence.  Thus
\[
 2P_z(S)=\cH(\delta S)+(1-\eta)b(S)+\rho(S).
\]
Substitute $b(S)=2F(S)-\cH(\delta S)$ from
Lemma~\ref{lem:cut-algebra}.  If $F(S)=0$, then
\[
 P_z(S)\le \frac{\eta n^2+n^{-20}}2<n^{-4}.
\]
If $F(S)\ge1$, the nonnegative $\eta\cH(\delta S)$ term and
$|\rho(S)|\le n^{-20}$ give
\[
 P_z(S)\ge1-\eta-n^{-20}/2>1/2.
\]
The finitely many smaller $n$ can be handled by learning all pairs, changing
only the absolute constant in the theorem.
\end{proof}

If the total multiplicity of an unordered pair is at most a fixed $C$, the
only change in this argument is
$\cH(\delta S)\le Cn^2$.  Thus the same choices
$\eta\le n^{-12}$ and $\tau=n^{-4}$ give the stated gap for all sufficiently
large $n$; the finitely many remaining sizes are again absorbed into the
absolute constant.

Define the implicit threshold digraph
\begin{equation}\label{eq:threshold-graph}
 J=\set{u\to v:c_{uv}>0\text{ and }p_{u\to v}\ge\tau},
 \qquad \tau=n^{-4}.
\end{equation}
The graph $J$ need not be a directed subgraph of $G$: it may orient an
underlying pair opposite to its sole input arc.  What is preserved is the
zero-outcut family.  Moreover, $J$ is not an orientation in the sense of
choosing exactly one direction per pair.  Since
$p_{u\to v}+p_{v\to u}=1$ and $\tau<1/2$, both ordered arcs may pass the
threshold.  The quantity $P_z(S)$ counts fractional mass once for each of
the $c_{uv}$ conceptual copies, whereas $J$ records only whether the
thresholded ordered pair exists; parallel copies do not create parallel arcs
in $J$.

\begin{lemma}[Threshold preservation]\label{lem:threshold}
For every $S\subseteq V$,
\[
 F(S)=0
 \quad\Longleftrightarrow\quad
 J\text{ has no arc from }S\text{ to }V\setminus S.
\]
Consequently, $G$ and $J$ have identical reachability relations.
\end{lemma}

\begin{proof}
If $F(S)=0$, Lemma~\ref{lem:gap} gives $P_z(S)<\tau$, so no individual
fractional arc leaving $S$ reaches threshold.  If $F(S)>0$, then
$P_z(S)>1/2$.  At most $Cn^2$ conceptual copies cross the cut (with
$C\le2$ for a simple digraph), so if all their leaving masses were below
$\tau=n^{-4}$ their sum would be at most $C/n^2<1/2$ for all sufficiently
large $n$, a contradiction.  As in Lemma~\ref{lem:gap}, the finitely many
remaining sizes are handled by adjusting the absolute constant.

For the reachability claim, $t$ is not reachable from $s$ exactly when the
set of vertices reachable from $s$ is an $s$-containing, $t$-excluding set
with no outgoing arc.  The family of such zero out-cuts is the same in $G$
and $J$.
\end{proof}

\begin{lemma}[Implicit search in $J$]\label{lem:implicit-bfs}
Given a pair-evaluable circuit for $z$ satisfying \eqref{eq:certificate},
reachability from $s$ in $J$ can be computed using $O(n\log n)$ additional
outgoing-cut queries.
\end{lemma}

\begin{proof}
We give the search and its invariant explicitly.  Maintain a discovered set
$R$, initially $\{s\}$, and a queue containing the discovered vertices that
have not yet been processed.  Both are determined by the public query
transcript.  When $u$ is removed from the queue, repeatedly form
\[
 T(u,R)=\{v\in V\setminus R:p_{u\to v}\ge\tau\}.
\]
This set is public even though the support of $H$ is hidden: membership uses
only $R$ and the pair-evaluable circuit for $z$, and does not test whether
$c_{uv}>0$.

If $T(u,R)$ is nonempty, use the rectangle identity to evaluate
\begin{equation}\label{eq:implicit-search-count}
 \cH(\{u\},T(u,R))
 =F(\{u\})+F(T(u,R))-F(T(u,R)\cup\{u\}).
\end{equation}
The value is positive exactly when some actual support pair $\{u,v\}$ has
$v\in T(u,R)$.  When it is positive, order $T(u,R)$ by the fixed public
vertex order and bisect it.  Apply
\eqref{eq:implicit-search-count} to the first half; if that count is
positive, retain the first half, and otherwise retain the second.  The
retained half always contains an actual neighbor, so after
$O(\log n)$ such tests a singleton $\{v\}$ remains.  Add $v$ to $R$ and to
the queue, recompute $T(u,R)$, and continue processing the same $u$.  When
$T(u,R)$ is empty or its count is zero, mark $u$ processed.  The algorithm
terminates when the queue is empty and returns $R$.

We prove that the returned set is exactly the set reachable from $s$ in
$J$.  Initially this is true for the only discovered vertex $s$.  Whenever
$v$ is discovered while processing $u$, the positive singleton count gives
$c_{uv}>0$, and membership in $T(u,R)$ gives
$p_{u\to v}\ge\tau$.  Hence $u\to v$ is an arc of $J$.  By induction, every
vertex placed in $R$ is reachable from $s$ in $J$.

Conversely, consider the final failed test for a processed vertex $u$.  At
that moment there is no actual pair from $u$ to any then-undiscovered
$v$ satisfying $p_{u\to v}\ge\tau$.  The discovered set only grows
afterwards, so there is also no $J$-arc from $u$ to the complement of the
final set $R$.  Every vertex in the final $R$ is eventually processed;
therefore $R$ is closed under outgoing arcs of $J$.  Since $s\in R$, no
directed path from $s$ can leave $R$, proving that no $J$-reachable vertex
was omitted.

There are at most $n-1$ successful searches, because each discovers a new
vertex, and at most one terminal failed search for each of at most $n$
processed vertices.  A successful or failed search uses
$O(\log n)$ count tests (an empty candidate set needs none), and each count
test uses three outgoing-cut queries, which may be issued in parallel.
Thus the total is $O(n\log n)$ queries.  The straightforward implementation
uses $O(n\log n)$ adaptive test rounds; the theorem charges only the number
of queries and does not claim a smaller round bound.
\end{proof}

\section{Weighted spectral sparsification from cut queries}
\label{sec:spectral}

The normal matrices arising later are weighted graph Laplacians whose support
is hidden but whose prospective conductance on a pair is endpoint-computable.
This section supplies the required oracle.

\begin{theorem}[Pair-evaluable weighted sparsification]
\label{thm:spectral}
Let $G$ be an unknown loopless directed multigraph on $[n]$, accessed through
the outgoing-cut oracle $F_G(S)=|E_G(S,[n]\setminus S)|$, where cardinality
counts arc copies.  Let $H$ be its
underlying undirected multigraph, and let
$c_e\in\{0,\ldots,C\}$ be the total number of arcs of $G$ in both directions
on the unordered pair $e$, where $C$ is a constant.  Let $w_e\ge0$ be
symmetric and pair-evaluable, and assume that nonzero weights can be compared
and output exactly.  For $0<\eps\le1/2$ and $0<\zeta<1/2$, there is a
randomized algorithm which uses
\[
 O\!\left(C\eps^{-2}n\log^8\!\frac{n}{\eps\zeta}\right)
\]
queries to $F_G$ and, with probability at least $1-\zeta$, outputs a
weighted graph $Q$ such that
\[
 (1-\eps)L\preceq L_Q\preceq(1+\eps)L,
 \qquad
 L=\sum_e c_ew_ea_ea_e^\top.
\]
The output has
$O(C\eps^{-2}n\log^6(n/(\eps\zeta)))$ unordered pairs.  The bound is
independent of the aspect ratio of the nonzero $w_e$.
\end{theorem}

\paragraph{Interface with the later normal systems.}
The exact-weight hypothesis in Theorem~\ref{thm:spectral} is an interface
condition, not an assertion that the analytic path state is stored exactly.
In the LP implementation, certified evaluation first freezes a positive
dyadic coefficient $\widehat D_e$ for every prospective hidden incidence
row.  Once the duplicate copies on a pair have a common row state, as proved
in Lemma~\ref{lem:block-symmetry}, their hidden normal-matrix contribution is
\[
        L_{\rm hid}
        =\sum_e c_e\widehat D_ea_ea_e^\top .
\]
Thus Theorem~\ref{thm:spectral} is called with
$w_e=\widehat D_e$.  This value is pair-evaluable because the row state is a
public endpoint circuit, and it is exact at this interface because
$\widehat D_e$ is dyadic.  The $O(n)$ explicit star-row contributions are
then added to the returned sparsifier offline.  In particular, constructing a
normal-system preconditioner or simultaneous leverage estimates never
requires learning whether a prospective pair is present.

No undirected-cut oracle is assumed.  The arc-by-arc proof of
Lemma~\ref{lem:cut-algebra} is linear in arc multiplicities and therefore
applies unchanged to this directed multigraph.  Thus, for $v\notin Y$, where
$Y$ is determined by the public transcript,
\[
 \sum_{u\in Y}c_{\{u,v\}}
 =F_G(\{v\})+F_G(Y)-F_G(Y\cup\{v\})
\]
by Lemma~\ref{lem:cut-algebra}.  Thus one filtered-neighborhood test costs
three outgoing-cut queries, and binary search for a qualifying neighbor costs
$O(\log n)$ such tests.  All bounds below include this translation.

\paragraph{Implicit residual supports.}
Every graph manipulated below is represented by a public survival predicate,
not by an edge list.  At recursion depth $r$, let $P_r(e)$ be a predicate
computable from the endpoints and the public transcript, and define
\[
 \mathcal E_r=\set{e:c_e>0\text{ and }P_r(e)},\qquad
 L_r=\sum_{e\in\mathcal E_r}c_ep_{\rm samp}^{-r}w_ea_ea_e^\top ,
 \qquad p_{\rm samp}=\tfrac14 .
\]
The symbol $p_{\rm samp}$ distinguishes this pair-sampling probability from
the cluster-sampling probability used inside one spanner call.
Initially $P_0(e)$ is the public test $w_e>0$.  When a public bundle
$B_r\subseteq\mathcal E_r$ has been found, draw a fresh public independent
Bernoulli hash $h_r(e)$ for every unordered endpoint pair, with
$\Prb[h_r(e)=1]=p_{\rm samp}$, and set
\begin{equation}\label{eq:survival-predicate}
 P_{r+1}(e)=P_r(e)\wedge[e\notin B_r]\wedge[h_r(e)=1].
\end{equation}
The hash is evaluated from the endpoint labels and a public seed, so
$P_{r+1}$ remains pair-evaluable.  Every filtered-neighbor query simply
intersects its candidate set with the current predicate.  Thus the algorithm
does not query whether an absent pair survived, and it materializes a pair
only when a filtered search actually returns it.  Conditional on the preceding
transcript, the hashes on the current actual pairs are independent fresh
quarter-sampling decisions.  This is the residual-support invariant used by
both the spanner construction and the recursion.  Within one spanner call,
the phase support $E_i$ is represented in the same way: conjoin $P_r$ with
the public block deletions and cluster predicates fixed by the completed
phases.  No predicate is changed while the queries for one frozen phase are
being answered.

\subsection{Filtered weighted spanners}

Give every support pair $e$ the distinct key
\[
        \operatorname{key}(e)=(1/w_e,\operatorname{ID}(e)),
\]
with zero-weight pairs excluded.  For a fixed vertex $v$, an arbitrary
public predicate $P(v,u)$ is enforced by placing in the query set only those
candidate endpoints satisfying $P$.  Corollary~\ref{cor:filtered-neighbor}
then supplies emptiness and minimum-neighbor queries.

We simulate the simultaneous weighted Baswana--Sen phases
~\cite{BaswanaSen2007}.  Let $E_i$ and $\mathcal C_i$ be the support and
clustering at the start of phase $i$, with $\mathcal C_1$ the singleton
clustering, and put $p=n^{-1/k}$.  In each phase $i<k$, sample every cluster
independently with probability $p$ and write $\mathcal S_i$ for the sampled
family.  A vertex in a sampled cluster makes no reports; set
$\mathcal R_i(v)=\varnothing$ for such a vertex.  For a vertex $v$ outside
the sampled clusters and
$C\in\mathcal C_i$, define
\[
 E_i(v,C)=\set{vu\in E_i:u\in C},\qquad
 e_i(v,C)=\arg\min_{e\in E_i(v,C)}\operatorname{key}(e),
\]
and let $\kappa_i(v,C)$ be the key of this edge, or $+\infty$ if the set is
empty.  The distinct pair IDs make every finite minimum unique.  Set
\[
 \sigma_i(v)=\min_{C\in\mathcal S_i}\kappa_i(v,C),
\]
and, when finite, let $C_i^\star(v)$ be its minimizing cluster.  The clusters
reported by $v$ are exactly
\[
 \mathcal R_i(v)=
 \begin{cases}
  \set{C:\kappa_i(v,C)<\sigma_i(v)}\cup\set{C_i^\star(v)},
       &\sigma_i(v)<+\infty,\\
  \set{C:\kappa_i(v,C)<+\infty},&\sigma_i(v)=+\infty.
 \end{cases}
\]
For every $C\in\mathcal R_i(v)$, insert $e_i(v,C)$ and record the entire
block $E_i(v,C)$ for deletion.  After all vertices have been processed, a
vertex with finite $\sigma_i(v)$ joins $C_i^\star(v)$; sampled clusters and
their members persist, while a vertex with infinite $\sigma_i(v)$ becomes
inactive.  Apply all recorded deletions simultaneously and then delete each
pair whose endpoints have the same new cluster label.  These operations
define $E_{i+1}$ and $\mathcal C_{i+1}$.  In phase $k$, take
$\mathcal S_k=\varnothing$ and report every neighboring cluster.

\paragraph{One phase, procedurally.}
The preceding definitions amount to the following frozen-snapshot procedure.
\begin{enumerate}[label=(\arabic*)]
\item Sample the old clusters, but do not yet change any cluster label or
      residual-support predicate.
\item For each vertex $v$ outside the sampled clusters, find its minimum-key
      edge to any sampled cluster; this determines $\sigma_i(v)$ and
      $C_i^\star(v)$.
\item Enumerate the old neighboring clusters whose minimum key is below
      $\sigma_i(v)$, together with $C_i^\star(v)$ when it exists.  For each
      reported cluster, insert its minimum-key edge and record the whole old
      vertex--cluster block for later deletion.
\item After every vertex has been examined against the same snapshot, apply
      the new cluster assignments, apply all recorded block deletions, and
      then delete pairs internal to a new cluster.
\end{enumerate}
All minimum, emptiness, and enumeration operations above are filtered searches:
their candidate endpoints satisfy the current public survival predicate and
the relevant public cluster predicate.

For example, suppose the old clusters adjacent to $v$ have minimum keys
\[
        \kappa_i(v,C_1)<\kappa_i(v,C_2)<\kappa_i(v,C_3),
\]
and among them only $C_2$ is sampled.  Then
$\sigma_i(v)=\kappa_i(v,C_2)$, the vertex reports $C_1$ and $C_2$, inserts
one minimum edge to each, records both corresponding blocks for deletion, and
joins $C_2$.  The block to $C_3$ survives this report unless it is removed by
another recorded deletion or becomes internal after the simultaneous
reclustering.  If $v$ already lies in a sampled cluster, it reports nothing.

Every quantity above refers to the one snapshot
$(E_i,\mathcal C_i,\mathcal S_i)$.  The filtered-neighbor oracle first finds
$\sigma_i(v)$, then enumerates precisely $\mathcal R_i(v)$, removing an
entire old cluster from the public candidate set after it is reported.
Delayed application of the block records therefore implements the displayed
phase exactly: one vertex's queries cannot change a later vertex's minimum,
report set, or cluster assignment.

\begin{lemma}[Query-efficient weighted spanner]
\label{lem:spanner}
For $k=\lceil\ln n\rceil$, a capped and restarted call returns a
$(2k-1)$-spanner containing $O(n\log n)$ pairs using
$n\log^{O(1)}(n/\zeta')$ outgoing-cut queries, except with probability
$\zeta'$.  Its query count is deterministically capped.
\end{lemma}

\begin{proof}
We verify stretch directly for these simultaneous phases.  Give pair $e$
length $\ell_e=1/w_e$, and let $\mathcal H_i$ be the pairs inserted before
phase $i$.  Inductively maintain the following invariant: for every
$e=uv\in E_i$ and each endpoint $x\in\{u,v\}$, if $x$ belongs to
$C\in\mathcal C_i$, then $\mathcal H_i$ contains an
$x$--$\operatorname{center}(C)$ path of at most $i-1$ edges, every one of
length at most $\ell_e$.  It is trivial for the singleton clustering.

Suppose $v$ joins $C_i^\star(v)$ through
$f=e_i(v,C_i^\star(v))$, and let $e=vu$ survive all recorded block
deletions, temporarily before the same-new-cluster deletion.  The old cluster
containing $u$ was not reported by $v$, so
\[
 \operatorname{key}(f)<\kappa_i(v,C_u)\le\operatorname{key}(e),
 \qquad\text{hence}\qquad \ell_f\le\ell_e.
\]
The old invariant for the endpoint of $f$ in $C_i^\star(v)$, followed by
$f$, gives the required path for $v$ with at most $i$ edges.  An endpoint
already in a sampled cluster keeps its old path.  A vertex that becomes
inactive reports every neighboring cluster, so no pair incident to it
survives.  This proves the claimed center-path property for every pair
surviving the block deletions; after pairs with equal new labels are removed,
it is precisely the invariant for phase $i+1$.

Now let $e=vu\in E_i(v,C)$ be deleted through a reported block, and put
$f=e_i(v,C)$.  Since $\ell_f\le\ell_e$, the invariant gives paths from $u$
and from the endpoint of $f$ in $C$ to the center of $C$, each using at most
$i-1$ edges of length at most $\ell_e$.  Together with $f$ they form a
$v$--$u$ path of length at most $(2i-1)\ell_e$.  If $e$ is instead deleted
because its endpoints receive the same new cluster label, the two center
paths use at most $i$ edges each, for total length at most $2i\ell_e$; this
can happen only when $i<k$.  In the final phase every remaining neighboring
cluster is reported, so every remaining pair obtains a path of length at
most $(2k-1)\ell_e$.  The output is therefore a $(2k-1)$-spanner.

For size and queries, condition on the frozen snapshot and order the old
clusters neighboring a fixed vertex by $\kappa_i(v,C)$.  If the vertex does
not itself lie in a sampled cluster, then in a sampling phase
$|\mathcal R_i(v)|$ is the position of the first sampled neighboring cluster,
truncated at the number of neighbors; if its own cluster is sampled, it
reports zero clusters.  Consequently
\[
 \Prb[|\mathcal R_i(v)|\ge j\mid E_i,\mathcal C_i]
 \le(1-p)^{j-1},\qquad
 \E[|\mathcal R_i(v)|\mid E_i,\mathcal C_i]\le p^{-1}.
\]
Moreover, $\E[|\mathcal C_{i+1}|\mid\mathcal C_i]=p|\mathcal C_i|$, so
$\E|\mathcal C_k|=np^{k-1}=n^{1/k}$.  The final phase reports at most
$n|\mathcal C_k|$ incidences.  If $R$ is the total number of reports, then
\[
        \E R\le(k-1)np^{-1}+n\E|\mathcal C_k|
             \le kn^{1+1/k}.
\]
At most one pair is inserted per report.  Abort once
$R>2kn^{1+1/k}$; Markov's inequality gives completion probability at least
$1/2$.  A trial performs one initial search per vertex per phase and
$O(1)$ binary searches per report, hence
$O((kn+R)\log n)$ filtered-neighborhood tests.  After
$\lceil\log_2(1/\zeta')\rceil$ independent capped trials, failure is at most
$\zeta'$, while the query count is deterministic.  With
$k=\lceil\ln n\rceil$, both the cap and output size are $O(n\log n)$.
\end{proof}

Multiplicity is ignored while building a spanner: all copies of a pair have
the same base length, and a selected pair is kept as one block.  Its
multiplicity is learned by a singleton instance of
Corollary~\ref{cor:filtered-neighbor}.

\subsection{Bundle leverage}

Fix one sparsification round.  Every current pair has conductance
$\omega_e^{(r)}=c_ep_{\rm samp}^{-r}w_e$, where
$p_{\rm samp}=1/4$ and $r$ is the number of previous pair-sampling rounds.
Construct $t$ spanners successively, deleting every selected
unordered pair before the next construction, and let $B$ be their union.

\begin{lemma}[Grouped bundle leverage]\label{lem:bundle}
Let each of the $t$ spanners have stretch $\alpha$.  If $e$ remains outside
$B$, then its leverage in the full current graph $G_r$ satisfies
\[
        \omega_e^{(r)}R_e[G_r]\le \frac{C\alpha}{t}.
\]
\end{lemma}

\begin{proof}
For every spanner $S_i$, the surviving pair $e=uv$ has a $u$--$v$ path
$P_i$ satisfying
\[
        \sum_{f\in P_i}\frac1{w_f}\le\frac{\alpha}{w_e}.
\]
The paths are pair-edge-disjoint because selected pairs are deleted between
successive spanners.  Send $1/t$ units of flow on every $P_i$.  The actual
resistance of $f$ is
$(c_fp_{\rm samp}^{-r}w_f)^{-1}\le p_{\rm samp}^r/w_f$, so the flow energy is at most
\[
 \frac1{t^2}\sum_{i=1}^t\sum_{f\in P_i}\frac{p_{\rm samp}^r}{w_f}
 \le \frac{p_{\rm samp}^r\alpha}{tw_e}.
\]
Thomson's principle bounds $R_e[G_r]$ by this energy.  Multiplication by
$\omega_e^{(r)}=c_ep_{\rm samp}^{-r}w_e$ and $c_e\le C$ proves the claim.  Internal vertices need
not be disjoint.
\end{proof}

\subsection{One sampling round}

Keep $B$ deterministically.  Independently retain every remaining unordered
pair with probability $p_{\rm samp}=1/4$, retaining all of its copies as one
block, and multiply a survivor's conductance by $1/p_{\rm samp}=4$.  Let $L$ be the Laplacian of the
\emph{full} graph at the start of the round, including $B$, and define on the
range of $L$
\[
 A_e=L^{\dagger/2}L_eL^{\dagger/2},
 \qquad L_e=\omega_e^{(r)}a_ea_e^\top,
\]
where $a_e=e_u-e_v$ is the canonical incidence vector.
Lemma~\ref{lem:bundle} gives $\norm{A_e}\le\lambda$, where
$\lambda=C\alpha/t$.  Moreover, the sum of $A_e$ over the sampled complement
is at most the identity projector.

If $\xi_e$ is the sampling indicator, the normalized zero-mean error is
\[
        Z_e=(\xi_e/p_{\rm samp}-1)A_e.
\]
For $p_{\rm samp}=1/4$,
\[
 \norm{Z_e}\le3\lambda,
 \qquad
 \E[(\xi_e/p_{\rm samp}-1)^2]=3,
 \qquad
 A_e^2\preceq\lambda A_e.
\]
Thus $\norm{\sum_e\E Z_e^2}\le3\lambda$.  Self-adjoint matrix Bernstein
~\cite{Tropp2012} yields the following.

\begin{lemma}[One-round concentration]\label{lem:one-round}
For $0<\delta\le1$,
\[
 \Prb\!\left[
  \norm{L^{\dagger/2}(\widehat L-L)L^{\dagger/2}}>\delta
 \right]
 \le 2n\exp\!\left(-\frac{\delta^2t}{8C\alpha}\right).
\]
\end{lemma}

The use of the full $L$ is important: the sampled complement alone may be
singular in directions controlled by the fixed bundle.

\subsection{Recursion and proof of Theorem~\ref{thm:spectral}}

Set
\[
 R=\left\lceil\log_4\frac{4n^2}{\zeta}\right\rceil,
 \quad
 \delta=\frac{\eps}{3R},
 \quad
 k=\lceil\ln n\rceil,
 \quad
 \alpha=2k-1,
\]
and choose
\[
 t=\Theta\!\left(C\alpha\delta^{-2}
                  \log\frac{nR}{\zeta}\right).
\]

\paragraph{Recursive procedure.}
Starting with $(r,P_r)=(0,P_0)$, for each
$r=0,\ldots,R-1$ perform the following steps.
\begin{enumerate}[label=(\arabic*)]
\item On the support $\mathcal E_r$, construct the $t$ successive spanners
      and let their explicit union be $B_r$.  When a pair first enters
      $B_r$, learn its multiplicity and store it with conductance
      $c_ep_{\rm samp}^{-r}w_e$.
\item Keep $B_r$ in the output, draw the fresh public hash $h_r$, and define
      the sampled residual by \eqref{eq:survival-predicate}.  Recurse on
      $(r+1,P_{r+1})$, where every surviving pair has conductance
      $c_ep_{\rm samp}^{-(r+1)}w_e$.
\item Return the union of $B_r$ and the graph returned by the recursive call.
\end{enumerate}
At depth $R$, use one filtered-neighbor emptiness test per vertex to verify
that $\mathcal E_R$ is empty.  If it is not, declare the capped attempt
unsuccessful; otherwise the union of the explicit bundles is the returned
graph.  The invariant above shows that every operation at level $r$ sees
exactly the Laplacian $L_r$, although its support has never been enumerated.

\paragraph{Approximation and termination.}
Conditioned on the transcript so far, the current hashes are
fresh, so Lemma~\ref{lem:one-round} applies.  If the recursive result
approximates the sampled residual within factors $a,b$, adding the fixed PSD
bundle and using the one-round event gives factors
$a(1-\delta),b(1+\delta)$ for the preceding graph.  Hence the final factors
are $(1-\delta)^R,(1+\delta)^R$, contained in
$[1-\eps,1+\eps]$.

The depth-$R$ check uses $n$ filtered-neighbor tests.  A fixed original pair
remains only if it avoids all
bundles and survives $R$ fresh quarter-samplings, so a union bound gives
probability at most $n^2 4^{-R}\le\zeta/4$.  Capped spanner-call failures and
the $R$ Bernstein events consume the remaining failure budget.

Finally,
\[
 Rt=O\!\left(C\eps^{-2}\log^5\frac{n}{\eps\zeta}\right).
\]
There are $Rt$ spanner calls, each costing
$O(n\log^3(n/(\eps\zeta)))$ queries after amplification.  This proves the
$O(C\eps^{-2}n\log^8(n/(\eps\zeta)))$ query bound.  Multiplying the capped
$O(n\log n)$ size per spanner by $Rt$ gives the asserted output size.

\section{An implicit rank-sensitive box-LP solver}
\label{sec:implicit-lp}

We next show how to run the required LP algorithm without enumerating hidden
rows.  Write a box LP in row-variable form as
\begin{equation}\label{eq:generic-box-lp}
 \min\{c_{\rm LP}^\top x:A^\top x=b_{\rm LP},\ \ell\le x\le u\},
 \qquad A\in\R^{M\times r}.
\end{equation}
For a supplied exactly feasible strict interior point $x^0$, satisfying
$A^\top x^0=b_{\rm LP}$, set
\[
 U=\max\set{
  \norm{1/(u-x^0)}_\infty,
  \norm{1/(x^0-\ell)}_\infty,
 \norm{u-\ell}_\infty,
 \norm{c_{\rm LP}}_\infty}.
\]

\subsection{IPM primer and interface dictionary}

This section is written so that the Lee--Sidford method can be used as an
interface theorem; familiarity with its interior-point analysis is not
required.  The following paragraph gives the geometric picture needed for the
rest of the proof.

For each row-coordinate $i$, let $\phi_i$ be the standard finite-box barrier
on $(\ell_i,u_i)$, and put
\[
 \Phi''(x)=\diag(\phi_i''(x_i)),\qquad W=\diag(w).
\]
At path parameter $t>0$ and positive path weights $w$, the method follows the
weighted barrier objective
\begin{equation}\label{eq:weighted-barrier-primer}
 f_t(x,w)=t c_{\rm LP}^{\top}x+\sum_{i=1}^M w_i\phi_i(x_i)
 \quad\text{subject to}\quad A^\top x=b_{\rm LP}.
\end{equation}
The exact minimizer of \eqref{eq:weighted-barrier-primer} is the
\emph{weighted center}.  A projected Newton step moves the current point
toward that center while remaining in the affine space.  The weights are not
fixed: Lee--Sidford continually chases a target vector built from Lewis
weights of a barrier-scaled copy of $A$.  This choice is what makes the number
of path steps depend on the column rank $r$, rather than on the possibly
quadratic number $M$ of row variables.

The mixed norm used to measure both centering error and weight motion is
\[
 \norm{v}_{w+\infty}
   \defeq \norm{v}_\infty+C_{\rm norm}\norm{v}_w,
 \qquad
 \norm{v}_w^2=v^\top Wv,
\]
where $C_{\rm norm}=\log^{O(1)}M$.  In this notation the centrality of a
strictly interior feasible point is
\begin{equation}\label{eq:centrality-primer}
 \delta_t(x,w)=
 \min_{\nu\in\R^r}
 \norm{
   \frac{t c_{\rm LP}+W\phi'(x)-A\nu}
        {\sqrt{W\phi''(x)}}
 }_{w+\infty}.
\end{equation}
Division and square roots are coordinatewise.  Thus $\delta_t(x,w)=0$ exactly
when the first-order stationarity equation holds for some equality multiplier
$\nu$; small centrality means that the current point is close enough to the
weighted center for the local Newton and path-change estimates to apply.
The stability result cited in Proposition~\ref{prop:ls-batches} allows the
normal solves and row reductions below to have inverse-polynomial error.

\paragraph{State-variable dictionary.}
The distinction between row space and constraint space is particularly
important here: $x,w,X$ have one coordinate per row, whereas $\nu,\pi$ and
normal-system right-hand sides have only $r$ coordinates.
\begin{center}
\small
\begin{tabular}{@{}p{.18\linewidth}p{.30\linewidth}p{.44\linewidth}@{}}
\toprule
Symbol & Role & Interpretation in this paper\\
\midrule
$M,r,A,b_{\rm LP}$ & dimensions, constraint matrix, and LP right-hand side
  & $A\in\R^{M\times r}$; hidden pair rows and public rows are row variables,
    while $A^\top x=b_{\rm LP}$ gives the $r$ equality constraints.\\
$x,\ell,u$ & primal point and box
  & $x$ is kept strictly between its coordinatewise endpoints during the
    path computation and remains exactly feasible.\\
$t,c_{\rm LP}$ & path parameter and objective
  & $t$ controls the objective term in \eqref{eq:weighted-barrier-primer};
    the execution has an artificial-cost phase and a true-cost phase.\\
$w,W,X$ & path weights
  & $W=\diag(w)$ and $X=\log w$ is the log-weight used by the chasing
    update; both are represented by public arithmetic circuits.\\
$\phi',\Phi''$ & barrier gradient and Hessian
  & $\Phi''=\diag(\phi_i''(x_i))$ supplies the local coordinate scaling.\\
$A_x,g(x)$ & Lewis target
  & $A_x=\Phi''(x)^{-1/2}A$ and $g(x)$ is the regularized Lewis-weight vector
    that the current $w$ chases.\\
$D,A^\top DA$ & primitive row scaling and normal matrix
  & $D$ changes from call to call; its hidden part is a weighted graph
    Laplacian and its public-row part is added explicitly.\\
$\pi,h$ & Newton multiplier and primal direction
  & solving for the $r$-vector $\pi$ produces an $M$-coordinate direction
    $h$ with $A^\top h=0$.\\
$\delta_t(x,w)$ & centrality
  & the mixed-norm stationarity residual in
    \eqref{eq:centrality-primer}.\\
\bottomrule
\end{tabular}
\end{center}

For the parameter choice used here, the regularized target weight is
\begin{equation}\label{eq:regularized-lewis-target}
 g_i(x)=c_0+w_{p_{\rm LS}}(A_x)_i,
 \qquad c_0=2r/M,\qquad
 p_{\rm LS}=1-1/\log(4M).
\end{equation}
Here $w_{p_{\rm LS}}(\mathcal A)$ denotes the vector of
$\ell_{p_{\rm LS}}$ Lewis weights of a matrix $\mathcal A$; it is a function
of the matrix and should not be confused with the current path-weight vector
$w$.  The value $c_0=2r/M$ is a constant-factor retuning of the fastest
parameter displayed by Lee--Sidford.  Their weight-function analysis permits
arbitrary nonnegative regularization~\cite[Theorem 29]{LeeSidford2020}, with
the constants below retuned accordingly.  We write
$k_{\rm LS}=O(\log M)$ for the corresponding regularized Lewis parameter in
the movement bounds.

\paragraph{Lee--Sidford-to-oracle crosswalk.}
The source algorithm is invoked only through the following interfaces.
\begin{center}
\small
\begin{tabular}{@{}p{.29\linewidth}p{.29\linewidth}p{.34\linewidth}@{}}
\toprule
Lee--Sidford operation & Object requested & Realization below\\
\midrule
coordinate update & one common formula for every row
  & a public pair-evaluable circuit, shared by duplicate copies\\
row reduction & sums, maxima, predicates, and mixed norms
  & bit-class aggregation and filtered-neighbor search\\
projected Newton step & solve $A^\top DA$
  & a spectral preconditioner followed by residual-certified Richardson
    iteration\\
Lewis update & simultaneous leverage estimates for a scaled $A_x$
  & a sparsifier inverse evaluated on every potential pair\\
weight chasing & support point of a mixed-norm ball
  & the one-dimensional KKT search in
    Subsection~\ref{subsec:mixed-support}\\
inexact oracle interface & approximate row reductions and normal solves
  & dyadic interface values from Lemma~\ref{lem:rounded-aggregate}, followed
    by the exact star repair in Section~\ref{sec:numerics}\\
\bottomrule
\end{tabular}
\end{center}

\paragraph{One representative primitive batch.}
Starting from the public state $(x,w,t)$, a centering-and-chasing batch
can be read as the following data flow.  First, common row formulas evaluate
the analytic barrier derivatives and requested diagonal row scaling.  Second,
certified evaluation freezes sufficiently accurate dyadic interface values, and row
reductions form the $r$-dimensional Newton right-hand side from those values.
Third, the hidden normal matrix for the frozen, pair-evaluable diagonal is
spectrally sparsified, and the resulting preconditioner is refined against the
full frozen matrix to obtain the Newton multiplier.  Fourth, when
the Lewis routine requests leverage scores, a more accurate sparsifier
inverse gives all of them simultaneously, after which the Lewis update is
coordinatewise.  Fifth, the mixed-norm support routine chooses the allowed
log-weight motion.  The primal direction is repaired on the explicit star,
and the batch-end domain check is made.  Not every batch invokes every
primitive, but Proposition~\ref{prop:ls-batches} asserts that this list is
exhaustive.  Section~\ref{sec:numerics} verifies that these approximations
meet the imported stable-solver interface.

\begin{proposition}[Lee--Sidford execution interface]
\label{prop:ls-batches}
If $A$ has full column rank and no zero row, the box-LP algorithm of Lee and
Sidford uses
\[
 O\!\left(\sqrt r\,\log^{13}M\,
          \log\frac{MU}{\eps}\right)
\]
primitive batches.  The execution in their Algorithms~1--3, including
\texttt{computeInitialWeight}, the artificial-cost phase, the cost switch,
and final centering, can be written so that a batch uses only a
polylogarithmic number of the following operations:
\begin{enumerate}[label=(\alph*)]
\item apply one common coordinate formula to every row, using that row's
      bounds, objective coefficient, current coordinate and weight, together
      with public scalars and $r$-dimensional vectors;
\item take sums, inner products, mixed norms, maxima, clipping decisions, and
      predicate-restricted reductions of row-coordinate data;
\item solve a positive definite system $A^\top DA$ for a positive diagonal
      $D$;
\item obtain simultaneous multiplicative leverage estimates for a diagonally
      scaled copy of $A$ and perform the coordinatewise Lewis iteration; and
\item return a support point of
      $\{v:\|v\|_\infty+C_{\rm norm}\|v\|_w\le R\}$.
\end{enumerate}
The proof of the source algorithm uses the leverage routine only through its
simultaneous multiplicative output guarantee.  In particular, its internal
random sketch is not part of the state-transition specification and may be
replaced by any routine with that guarantee.  With exact versions of
(a)--(e), the execution returns an $\eps$-suboptimal point in the box.
The same conclusion, with constant-factor retuning of the invariant margins,
holds when row reductions, mixed-ball support points, and normal solves meet
the inverse-polynomial additive, multiplicative, and energy-norm accuracies
requested by the execution.
\end{proposition}

\begin{proof}
Theorem~1 and Algorithm~3 of the arXiv version of
Lee--Sidford~\cite{LeeSidford2020} give the displayed rank-sensitive count and
spell out the artificial-cost and true-cost phases.  Their Theorem~40 reduces
each path step to a centering call, a mixed-ball projection, and approximate
Lewis weights.  Appendix~B implements the initial and subsequent Lewis
weights by coordinatewise arithmetic and simultaneous leverage estimates;
its Theorem~39 assumes only the resulting multiplicative estimates.
Appendix~D, culminating in Theorem~62, is precisely the support problem in
(e) after the diagonal change of variables used below.  The remaining lines
of Algorithms~1--3 are the coordinate maps and reductions in (a)--(b).
Grouping the polylogarithmically many occurrences in one path step and
absorbing that factor into $\log^{13}M$ gives the stated batches.  This
enumeration includes initialization and the cost switch; no additional
row-access operation is used by the source execution.

For the last sentence of the proposition, the journal version explicitly
uses approximate Lewis weights and notes that inverse-polynomially accurate
linear solves suffice.  The detailed stability argument appears in
Section~8, especially Theorem~28, of the earlier primal path-following
paper~\cite{LeeSidfordPathII}.  That theorem is an interface statement in
energy-relative solve error; decreasing all requested accuracies by a fixed
polynomial factor only adds logarithmically many refinement iterations.
\end{proof}

\begin{lemma}[Conditional Lewis computation]\label{lem:conditional-lewis}
Fix a matrix $\mathcal A$ with full column rank.  In the approximate- and
initial-weight routines of Lee--Sidford, replace every randomized
leverage-score call by arbitrary simultaneously valid
multiplicative estimates of the requested accuracy.  Conditional on those
estimates, all subsequent updates are deterministic coordinatewise formulas,
and the conclusions and iteration bounds of Theorems~39 and~58, using
Lemma~60, of
Lee--Sidford~\cite{LeeSidford2020} continue to hold.  In particular,
the initial-weight routine may be started at $p=2$ with simultaneous
leverage estimates for $\mathcal A$.  At the $\log^{-O(1)}M$ accuracies requested by
the path algorithm, every accuracy-dependent factor is absorbed by the number
of leverage calls charged in Proposition~\ref{prop:ls-batches}.
\end{lemma}

\begin{proof}
For $p=2$ the Lewis weights are exactly the leverage scores of $\mathcal A$.
Algorithm~7 of the cited paper then changes $p$ by its displayed homotopy
schedule and invokes \texttt{computeApxWeight}.  Once the simultaneous
leverage inequalities requested by a call are supplied, Algorithms~6--7 use
only coordinatewise powers, medians, clipping, and arithmetic.  The proofs of
Theorems~39 and~58 and Lemma~60 condition only on those simultaneous
inequalities; the
Johnson--Lindenstrauss signs are one implementation of the interface, not an
additional hypothesis about the update.  Induction over the homotopy calls
therefore proves the statement.  The final call in Algorithm~7 accepts the
requested accuracy parameter.  Its polynomial dependence on reciprocal
accuracy is polylogarithmic for the accuracies used here and is already
included in the $\log^{13}M$ batch bound.  Higher fixed-point precision at an
interface is not being conflated with a more accurate Lewis iteration.
\end{proof}

Section~\ref{sec:numerics} checks that our cut-query implementation meets
this inexact interface while preserving affine feasibility exactly.

\begin{theorem}[Implicit incidence box LP]\label{thm:implicit-lp}
Fix a constant $C$.  Consider \eqref{eq:generic-box-lp}, whose $r$ columns
are indexed by $[r]$, an accuracy $0<\eps<1$, and a supplied point $x^0$
satisfying
\[ 
        A^\top x^0=b_{\rm LP},\qquad \ell<x^0<u.
\]
The hidden support is the symmetrization of a directed multigraph accessed by
its outgoing-cut oracle.  Suppose:
\begin{enumerate}[label=(\roman*)]
\item for each unordered pair $\{u,v\}\subseteq[r]$, the hidden rows consist
      of $c_{uv}\le C$ copies of $e_u-e_v$;
\item all copies on a pair have identical pair-evaluable interval bounds,
      objective coefficient, and initial coordinate;
\item the remaining $O(r)$ rows are explicit, all rows have Euclidean norm at
      most $\sqrt2$, and the explicit rows include a finite-box $+I_r$ block;
      and
\item every explicit matrix coefficient, right-hand side, box endpoint,
      objective coefficient, and coordinate of $x^0$ is dyadic, has
      $O(\log(MU/\eps))$ total bit length, and has absolute value
      $(MU/\eps)^{O(1)}$; the hidden incidence coefficients are integral; and
\item $M,U,\eps^{-1}=r^{O(1)}$.
\end{enumerate}
Then, with probability at least $1-r^{-\Omega(1)}$, the algorithm returns a
pair-evaluable represented point $x$ satisfying
\[
 \ell\le x\le u,\qquad A^\top x=b_{\rm LP},\qquad
 c_{\rm LP}^\top x\le\OPT+\eps,
\]
using $r^{3/2}\log^{O(1)}r$ outgoing-cut queries.  Every represented primal
iterate satisfies its equality constraints exactly.
\end{theorem}

For the Phase-I instance, $b_{\rm LP}=d$, $M\le n^2+2n$, and $r=n$.  Moreover
$|b_v|\le n-1$, hence $K_0\le n$; the artificial start slack is at least
$(K_0+1)/2$, the hidden start slack is constant, and the above width
parameter satisfies $U\le2n$.  Thus the theorem costs
$n^{3/2}\log^{O(1)}n$ queries.

\paragraph{Proof boundary and dependency roadmap.}
The proof of Theorem~\ref{thm:implicit-lp} begins here and ends at the close of
Section~\ref{sec:numerics}.  Two facts are imported from Lee--Sidford:
Proposition~\ref{prop:ls-batches} gives the exact-arithmetic batch count and
the exhaustive primitive interface, while
Lemma~\ref{lem:conditional-lewis} records that their Lewis-weight analysis
needs only simultaneous leverage guarantees.  The remainder of the proof has
four internal components:
\begin{enumerate}[label=(\arabic*)]
\item Lemma~\ref{lem:block-symmetry} collapses identical conceptual copies to
      one pair circuit without changing their linear contribution to a normal
      matrix.
\item Lemma~\ref{lem:transcript-simulation} inducts over the public transcript
      and reduces coordinate operations and row aggregates to cut queries.
\item The next two subsections implement the non-coordinate primitives:
      normal solves and leverage scores via Theorem~\ref{thm:spectral}, and
      mixed-norm support via a scalar KKT search.
\item Section~\ref{sec:numerics} checks the imported inexact interface and
      uses the explicit star rows to preserve affine feasibility exactly.
\end{enumerate}
Each batch then costs $r\log^{O(1)}r$ queries.  Multiplying by the
$\sqrt r\log^{O(1)}r$ batch count proves the claimed query bound; the same
induction shows that no operation tests an unqueried adjacency bit.

\subsection{Duplicate-copy symmetry}

A conceptual copy is a separate row coordinate, not a row whose incidence
vector has been multiplied by the pair multiplicity.  In general, if the
rows of $A$ are $a_i^\top$, then
$A^\top DA=\sum_iD_i a_i a_i^\top$.  Hence $c_{uv}$ identical rows, each with
diagonal value $D_{uv}$, contribute
$c_{uv}D_{uv}a_{uv}a_{uv}^\top$.  Replacing them by the single scaled row
$c_{uv}a_{uv}^\top$ would incorrectly square the multiplicity.  The following
lemma justifies retaining only one represented formula while still using the
former, linear contribution.

\begin{lemma}[Block symmetry]\label{lem:block-symmetry}
If all hidden copies of an unordered pair start identically, every operation
of the compiled box-LP algorithm preserves their equality.  Consequently,
every hidden normal matrix is
\begin{equation}\label{eq:hidden-normal}
        L_D=\sum_{u<v}c_{uv}D_{uv}a_{uv}a_{uv}^\top,
\end{equation}
not a matrix with coefficient $c_{uv}^2D_{uv}$.
\end{lemma}

\begin{proof}
Equal copies have the same row, box, objective coefficient, coordinate,
barrier derivatives, and weight.  A Newton coordinate is a common
coordinatewise expression plus the endpoint potential difference
$a_{uv}^\top\pi$.  Its leverage is
$D_{uv}a_{uv}^\top(A^\top DA)^{-1}a_{uv}$, also common.  Lewis updates use
coordinatewise powers, affine combinations, medians, and clipping.  The
mixed-ball optimizer below assigns equal outputs to equal data.  Induction
therefore proves equality throughout.  Summing the $c_{uv}$ identical
rank-one contributions gives \eqref{eq:hidden-normal}.
\end{proof}

Independent JL signs cannot be assigned to indistinguishable copies through
the cut oracle.  We avoid them by obtaining deterministic simultaneous
leverage estimates from a sparsifier inverse.

\begin{lemma}[Transcript simulation invariant]\label{lem:transcript-simulation}
At the beginning and end of every batch in
Proposition~\ref{prop:ls-batches}, each hidden row coordinate is represented
by one public pair-evaluable circuit, common to all copies of its unordered
pair.  Every operation (a)--(e) can be simulated with
$r\log^{O(1)}r$ outgoing-cut queries, and its represented outputs again satisfy
this invariant.
\end{lemma}

\begin{proof}
Initially the bounds, objective, and primal coordinate are public pair
circuits by assumption.  There is no circular appeal to an already computed
initial weight.  At $p=2$, the Lewis weights of
$\mathcal A_0=\Phi''(x^0)^{-1/2}A$ are its leverage scores.  One simultaneous-leverage
call therefore produces a pair-evaluable approximation to the
regularized $p=2$ target.  Lemma~\ref{lem:conditional-lewis} then inducts
through the $p=2$-to-$p_{\rm LS}$ homotopy, so its final log-weight is also a
public pair circuit.

This first leverage call is licensed directly from the supplied start, whose
slack is at least $1/U$, so
$\Phi''(x^0)$ has polynomial upper and lower bounds, while the explicit
$+I_r$ block gives a polynomial lower spectral bound for its normal matrix.
To make the induction quantitative, let $w_p(\mathcal A_0)$ be the exact
$p$-Lewis
weights for $p\in[p_{\rm LS},2]$, let $j$ minimize their coordinates, put
$a=1-2/p\le0$, and write $\alpha_i^\top$ for a row of $\mathcal A_0$.  The Lewis fixed-point
identity and
\[
 \mathcal A_0^\top W_p^a\mathcal A_0\preceq
 (w_p)_j^a \mathcal A_0^\top\mathcal A_0
\]
give
\[
 (w_p)_j=(w_p)_j^a \alpha_j^\top
 (\mathcal A_0^\top W_p^a\mathcal A_0)^{-1}\alpha_j
 \ge \alpha_j^\top
       (\mathcal A_0^\top\mathcal A_0)^{-1}\alpha_j
 \ge(MU)^{-O(1)}.
\]
The last inequality uses the polynomial spectral upper bound, the absence of
zero rows, and the dyadic input-bit hypothesis; a nonzero explicit
coefficient and every scaled hidden incidence coefficient have
inverse-polynomial magnitude.
Every exact Lewis weight is at most one.  During the initial-weight homotopy
$x=x^0$ is fixed, and Lemma~60 together with the median clipping in
Algorithm~7 therefore keeps every intermediate approximate scaling between
inverse polynomials and polynomials.  Each successive scaled normal matrix is
thus polynomially conditioned, so the same direct argument licenses every
call counted by Proposition~\ref{prop:ls-batches}.

A common coordinate formula preserves the pair-circuit property.  Sums,
inner products, incidence products, and predicate-restricted sums use
Lemmas~\ref{lem:aggregate} and~\ref{lem:rounded-aggregate}; mixed norms add
one square root after such sums.  Maxima, $\ell_\infty$ norms, emptiness
tests, and clipping predicates use Lemma~\ref{lem:hidden-max}.  When a
certified interval straddles a threshold, either branch changes the
corresponding coordinate by at most the interval width.  We choose that width
below the inverse-polynomial tolerance in
Proposition~\ref{prop:ls-batches}.

Normal solves and simultaneous leverage estimates are implemented in the
next subsection, and the mixed-ball support point in the following one.  Their
outputs consist of public $r$-vectors, public scalars, or a common
pair-evaluable row formula, so they also preserve the invariant.  All random
hashes used by the spectral calls are public transcript data.

For completeness, the artificial cost is kept as the exact elementary
circuit
\[
        (c_{\rm art})_i=-\exp(X_i^0)\phi_i'(x_i^0).
\]
Consequently $x^0$ is the exact center of the analytic artificial objective;
only evaluations of this fixed circuit are rounded at aggregate interfaces.
The cost switch replaces this circuit by the supplied true-cost circuit and
does not inspect hidden support.  This covers the initialization, both path
phases, and the final centering listed in
Proposition~\ref{prop:ls-batches}.
\end{proof}

\subsection{Normal systems and leverage scores}

\begin{lemma}[Certified dyadic inverse interface]
\label{lem:dyadic-inverse}
Let $K$ be an explicit dyadic positive definite $r\times r$ matrix with
\[
        K_-I\preceq K\preceq K_+I,
\]
where $K_+,$ $K_-^{-1}$, and $\vartheta^{-1}$ are polynomially bounded in
$MU/\eps$.  Offline certified arithmetic can output a dyadic symmetric $Q$
with $O(\log(MU/\eps))$ bits per entry such that
\[
 (1-\vartheta)K^{-1}\preceq Q
 \preceq(1+\vartheta)K^{-1}.
\]
\end{lemma}

\begin{proof}
Compute certified intervals for the entries of $K^{-1}$ until their widths
are at most
$\delta=\vartheta/(4rK_+) $, and take the symmetric dyadic midpoints.
Certified Gaussian elimination or interval Cholesky terminates because
$K\succeq K_-I$.  For $E=Q-K^{-1}$, the entrywise bound gives
$\norm E_2\le r\delta\le\vartheta/(2K_+)$ after harmless constant
rounding.  Since $K^{-1}\succeq K_+^{-1}I$, this implies
$-\vartheta K^{-1}\preceq E\preceq\vartheta K^{-1}$.
The magnitudes of inverse entries are at most $K_-^{-1}$, while the requested
absolute accuracy is inverse polynomial, proving the bit bound.
\end{proof}

\paragraph{How a normal matrix enters Theorem~\ref{thm:spectral}.}
This is the precise bridge between the IPM and the graph primitive.  During
one normal-system or leverage call, the analytic positive row scaling $D$ is
held fixed.  Certified interval evaluation chooses a positive dyadic value
$\widehat D_{uv}$ for every potential hidden pair, using the same value for
all conceptual copies of that pair.  The resulting hidden contribution is
\[
 \widehat L_{\rm hid}
   =\sum_{u<v}c_{uv}\widehat D_{uv}
      a_{uv}a_{uv}^\top .
\]
Thus Theorem~\ref{thm:spectral} is instantiated with the exact
pair-evaluable edge weight
\[
        w_{\{u,v\}}^{\rm spec}=\widehat D_{uv}.
\]
The public dyadic circuit can compare and output these weights exactly, as
required by that theorem; the unknown multiplicity $c_{uv}$ remains part of
the hidden graph and is learned only through cut queries.  In particular, the
spectral routine never receives an unevaluated exponential or barrier
derivative.  All explicit-row contributions are instead formed exactly and
added to the returned hidden sparsifier.  Consequently the graph theorem is
applied to precisely the hidden summand of the frozen dyadic normal matrix,
not directly to the analytic matrix and not to a graph in which duplicate
rows have been combined by scaling their incidence vectors.

For an analytic system $L=A^\top DA$, certified evaluation first forms a
positive dyadic diagonal $\widehat D$ satisfying
\begin{equation}\label{eq:dyadic-diagonal}
        (1-\beta_N)D\preceq\widehat D\preceq(1+\beta_N)D,
        \qquad \beta_N=(MU/\eps)^{-C_N}.
\end{equation}
Corollary~68 and Lemma~69 of Lee--Sidford~\cite{LeeSidford2020} give a
polynomial envelope for every normal system used by their execution.  A
sufficiently large absolute $C_N$ therefore makes this possible with
logarithmic precision.
These dyadic values are frozen for the entire spectral call and all
subsequent refinement steps belonging to that normal solve.  Put
$\widehat L=A^\top\widehat DA$.  Apply
Theorem~\ref{thm:spectral} with constant accuracy to the hidden contribution
of $\widehat L$ and add every explicit-row contribution exactly.  PSD addition
gives an explicit preconditioner $P$ with
\[
        \tfrac12\widehat L\preceq P\preceq2\widehat L .
\]
A full product by $\widehat L$ is computed exactly for the dyadic interface
values by Lemma~\ref{lem:aggregate}.  Preconditioned Richardson therefore
converges geometrically for the frozen dyadic system.  Here and below an
explicit positive definite preconditioner is not inverted as an exact
rational matrix.  Lemma~\ref{lem:dyadic-inverse} instead produces a
dyadic matrix $Q_P$ with
\begin{equation}\label{eq:dyadic-preconditioner-inverse}
 (1-\beta_P)P^{-1}\preceq Q_P\preceq(1+\beta_P)P^{-1},
 \qquad \beta_P\le1/20 .
\end{equation}
The polynomial spectral envelope supplies the hypotheses of the lemma.
Preconditioned Richardson contracts the energy error by a constant factor,
so $O(\log(MU/\eps))$ iterations give the inverse-polynomial relative
accuracy requested in Proposition~\ref{prop:ls-batches}.  Each iteration
multiplies by the full frozen matrix using
Lemma~\ref{lem:rounded-aggregate}; hence one such solve uses
$r\log^{O(1)}r$ queries.  Section~\ref{sec:numerics} compares the resulting
direction with the analytic one and applies the exact star repair.

For leverage scores, request a $(1\pm\theta)$ sparsifier $L_S$ of the full
dyadic normal matrix $\widehat L$, with
$\theta=1/\polylog r$.  Lemma~\ref{lem:dyadic-inverse} produces a
dyadic symmetric matrix $Q_S$ satisfying
\begin{equation}\label{eq:dyadic-sparsifier-inverse}
 (1-\theta/4)L_S^{-1}\preceq Q_S
 \preceq(1+\theta/4)L_S^{-1}.
\end{equation}
For a
hidden copy on $e=uv$, define
\begin{equation}\label{eq:leverage-estimate}
 \widehat\sigma_e
   =\widehat D_e\bigl((Q_S)_{uu}+(Q_S)_{vv}
                       -2(Q_S)_{uv}\bigr).
\end{equation}
Inverse monotonicity and \eqref{eq:dyadic-sparsifier-inverse} give
\[
 (1-3\theta)\widehat D_ea_e^\top\widehat L^{-1}a_e
 \le\widehat\sigma_e\le
 (1+3\theta)\widehat D_ea_e^\top\widehat L^{-1}a_e
\]
for $\theta$ below an absolute constant.
Together with \eqref{eq:dyadic-diagonal}, this is a
$1\pm O(\theta+\beta_N)$ estimate of the analytic leverage score.  Requesting
half the tolerance from each source gives the tolerance required by
Lee--Sidford.  The conclusion holds simultaneously for every actual or
potential pair.  Formula~\eqref{eq:leverage-estimate} is pair-evaluable and
common across duplicate copies, so the approximate Lewis updates remain
coordinatewise pair circuits.

\subsection{The mixed-norm support primitive}
\label{subsec:mixed-support}

For a symmetric compact convex set $\mathcal K$, write
\[
 h_{\mathcal K}(p)\defeq\max_{v\in\mathcal K}\ip{p}{v},
 \qquad \norm{p}_{\mathcal K}\defeq h_{\mathcal K}(p).
\]
The second notation is only shorthand for the support function.  The chasing
update requires a support point of
\[
 \mathcal U_R=
 \set{u:\norm{u}_\infty+C_{\rm norm}\norm{u}_w\le R},
 \qquad \norm{u}_w^2=u^\top Wu.
\]
This primitive is the decision step in Lewis-weight chasing.  Given the
current gradient-like objective $p$, it chooses the admissible log-weight
motion with nearly maximum predicted progress.  The $\ell_\infty$ term
prevents any one row weight from moving too far, while the weighted
$\ell_2$ term controls aggregate motion.  Because the coordinates include
hidden rows, the optimizer cannot be constructed by scanning the support.

\paragraph{Algorithmic view.}
The computation has four steps.  First, diagonally rescale the ball to
\eqref{eq:normalized-mixed-support}.  Second, introduce a scalar
$\tau\in[0,1]$ that splits the unit budget between the Euclidean radius
$1-\tau$ and the coordinate caps $\tau\kappa_i$.  For a fixed $\tau$, the KKT
conditions reduce the optimizer to one scalar multiplier $\lambda$, found by
bisection from predicate-restricted aggregates.  Third, maximize the resulting
one-dimensional concave value function $G(\tau)$ by certified interval
search.  Finally, undo the diagonal scaling and, if necessary, rescale
the dyadic output conservatively to preserve feasibility.  The formulas
below establish the relative support guarantee and the
query bound for these four steps.  The letter $\tau$ here is only this local
budget split and is unrelated to the central-path parameter $t$.

If $R=0$, the routine returns zero exactly.  Assume $R>0$.  For an objective
vector $p$, set
\[
 x=\frac{C_{\rm norm}}R W^{1/2}u,
 \qquad \kappa_i=C_{\rm norm}\sqrt{w_i},
 \qquad a=\frac{R}{C_{\rm norm}}W^{-1/2}p.
\]
Then $\ip{p}{u}=\ip{a}{x}$ and the normalized problem is
\begin{equation}\label{eq:normalized-mixed-support}
 \max a^\top x
 \quad\text{s.t.}\quad
 \norm{x}_2+\norm{\kappa^{-1}x}_\infty\le1.
\end{equation}
For $\tau\in[0,1]$, define
\begin{equation}\label{eq:mixed-value-function}
 G(\tau)=\max\set{a^\top x:\norm{x}_2\le1-\tau,
                         \ |x_i|\le \tau\kappa_i\ \text{for every }i}.
\end{equation}
The endpoints have the unique feasible optimizer $x=0$.  If
$0<\tau<1$ and
\[
 \sum_{i:a_i\ne0}\tau^2\kappa_i^2\le(1-\tau)^2,
\]
the optimizer is
$x_i=\operatorname{sign}(a_i)\tau\kappa_i$ for $a_i\ne0$ and
$x_i=0$ otherwise.  In the remaining case, the KKT conditions give
\begin{equation}\label{eq:mixed-kkt}
 x_i=\operatorname{sign}(a_i)
       \min\set{\tau\kappa_i,\frac{|a_i|}{2\lambda}},
\end{equation}
where the unique $\lambda>0$ satisfies
\begin{equation}\label{eq:mixed-lambda}
 \sum_i\min\set{\tau^2\kappa_i^2,\frac{a_i^2}{4\lambda^2}}
       =(1-\tau)^2.
\end{equation}
Every predicate and sum in \eqref{eq:mixed-lambda}, and the resulting
objective value, is a predicate-restricted pair aggregate.

We record the relative guarantee needed later.  The function $G$ is concave:
interpolating feasible solutions at $\sigma,\tau$ gives a feasible solution at the
same interpolation of the parameters.  It is also Lipschitz, with
\begin{equation}\label{eq:mixed-value-lipschitz}
 |G(\tau)-G(\sigma)|\le L_a|\tau-\sigma|,
 \qquad L_a=\max\set{\norm{a}_2,\sum_i|a_i|\kappa_i}.
\end{equation}
Indeed, for $\sigma<\tau$, scale an optimizer at $\sigma$ by
$(1-\tau)/(1-\sigma)$ to obtain a feasible point at $\tau$, losing at most
$\norm{a}_2(\tau-\sigma)$; scaling an optimizer at $\tau$ by
$\sigma/\tau$ gives a feasible point at $\sigma$, losing at most
$(\tau-\sigma)\sum_i|a_i|\kappa_i$.

If $a=0$, return zero.  Otherwise the one-coordinate point
\[
 x_i=\operatorname{sign}(a_i)\frac{\kappa_i}{1+\kappa_i},
 \qquad x_j=0\quad(j\ne i),
\]
is feasible for every $i$ with $a_i\ne0$.  Hence, if
$H=\max_{\tau}G(\tau)=h_{\mathcal U_R}(p)$, then
\begin{equation}\label{eq:mixed-support-lower-bound}
 H\ge\nu_0\defeq
 \max_{i:a_i\ne0}\frac{|a_i|\kappa_i}{1+\kappa_i}.
\end{equation}
At an implemented call the objective is frozen on an inverse-polynomial
dyadic mesh.  If its support value is below the requested additive tolerance,
the zero vector is already an admissible answer.  Otherwise a nonzero frozen
coordinate has inverse-polynomial magnitude.  The source weight bounds
(Lee--Sidford~\cite[Lemma~67]{LeeSidford2020}) make
$W^{\pm1/2}$ polynomially bounded, and every positive radius requested by
their path schedule is inverse polynomial in $MU/\eps$.  Thus
$\nu_0\ge(MU/\eps)^{-O(1)}$ in the only case where a relative search is
needed.

Bisection on \eqref{eq:mixed-lambda}, followed if needed by a conservative
common rescaling, evaluates $G(\tau)$ and returns a feasible sign-aligned point
to any additive accuracy $\zeta$.  A noise-robust one-dimensional concave
search then maximizes $G$: compare certified value intervals at two interior
trisection points, discard a side when the intervals separate, and otherwise
retain the better queried point.  Concavity and
\eqref{eq:mixed-value-lipschitz} show that
$O(\log(L_a/\zeta))$ evaluations achieve additive error $O(\zeta)$.
Taking $\zeta\le\gamma\nu_0/32$ and accounting for both searches yields a
feasible, sign-aligned $u\in\mathcal U_R$ satisfying
\begin{equation}\label{eq:relative-mixed-support}
        \ip{p}{u}\ge(1-\gamma)h_{\mathcal U_R}(p).
\end{equation}
For inverse-polynomial $\gamma$, all scalar searches use $O(\log(MU/\eps))$
aggregate evaluations, and the query cost is $r\log^{O(1)}r$.
All sums involving the analytic weights are evaluated by certified dyadic
intervals.  The routine rescales using an upper interval for the mixed norm,
so the returned point is feasible for the \emph{analytic} ball.
Lemma~\ref{lem:rounded-aggregate} and the Lipschitz
bound \eqref{eq:mixed-value-lipschitz} show that width $\beta$ in each
interface changes the support value by at most
$\poly(MU/\eps)\beta$.  Choosing $\beta$ below the source tolerance includes
this loss in $\zeta$.

The list in Proposition~\ref{prop:ls-batches} is exhaustive.
Lemma~\ref{lem:transcript-simulation} handles its coordinate operations and
reductions, while the two preceding subsections handle its remaining
primitives.  Consequently no step branches on an unqueried adjacency bit.
This completes the cut-query simulation part of the proof of
Theorem~\ref{thm:implicit-lp}.  It remains only to check the stable-solver
accuracy interface and exact affine feasibility; that is the purpose of the
next section.

\section{The solver interface and exact feasibility}
\label{sec:numerics}

The numerical layer is short once the right boundary is drawn.  We use the
stable Lee--Sidford execution as an imported solver theorem and prove only
that our cut-query primitives meet its accuracy requirements.  This is the
one place where the explicit $+I_r$ block is used algorithmically.

At a current point let
\[
 H=W\Phi''(x),\qquad D=H^{-1},\qquad
 L=A^\top DA,\qquad q=t c_{\rm LP}+W\phi'(x).
\]
The exact projected Newton direction has the form
\begin{equation}\label{eq:exact-newton}
 h=-D(q-A\pi),\qquad L\pi=A^\top Dq,\qquad A^\top h=0.
\end{equation}
Let $R_\star\in\mathbb R^{M\times r}$ embed a vector on the explicit
$+I_r$ rows.  Then
\begin{equation}\label{eq:star-right-inverse}
             A^\top R_\star=I_r.
\end{equation}

\begin{lemma}[Cut-query solver interface]\label{lem:solver-interface}
Under the hypotheses of Theorem~\ref{thm:implicit-lp}, there is an absolute
constant $C_{\rm stab}$ for which the following implementation is valid.
Evaluate every row reduction and mixed-norm support call to additive or
multiplicative error at most
\[
             \alpha=(MU/\eps)^{-C_{\rm stab}},
\]
and solve every normal system to relative $\alpha$ error in its energy norm.
After forming the resulting raw primal direction $\widetilde h$, replace it
by
\begin{equation}\label{eq:star-repair}
 h_{\rm rep}=\widetilde h-R_\star(A^\top\widetilde h).
\end{equation}
Then the execution satisfies the inexact interface of
Proposition~\ref{prop:ls-batches}, every represented primal iterate obeys
$A^\top x=b_{\rm LP}$ exactly, and the final point lies in the box and has
objective at most $\OPT+\eps$.

Each primitive batch uses $r\log^{O(1)}r$ cut queries.  All dyadic values
passed to bit-class aggregation need only
$\log^{O(1)}(MU/\eps)$ bits.
\end{lemma}

\begin{proof}
We first record what is imported.  The inexact path-following analysis of
Lee--Sidford~\cite[Section~8, especially Theorem~28]{LeeSidfordPathII}
shows that inverse-polynomial energy-relative errors in the linear solves can
be absorbed by constant-factor slack in the centering invariants.  The later
rank-sensitive formulation uses approximate Lewis weights explicitly
~\cite[Theorems~39--40]{LeeSidford2020}.  Its Corollary~68 bounds every
barrier Hessian encountered by a polynomial in $MU/\eps$, and Lemma~69 gives
the corresponding polynomial relative conditioning of the normal systems.
Lemma~67 bounds path weights between inverse polynomials and constants; the
finite-box barrier formulas and the width parameter $U$ give the complementary
lower Hessian bounds.
Consequently there is one fixed polynomial accuracy that suffices
simultaneously for all the row reductions, support computations, leverage
estimates, and normal solves listed in
Proposition~\ref{prop:ls-batches}.  We choose $C_{\rm stab}$ larger than its
degree.

The coordinate reductions meet this accuracy directly.  Certified evaluation
freezes a row circuit to dyadic width $\beta$.  By
Lemma~\ref{lem:rounded-aggregate}, every scalar aggregate changes by at most
$M\beta$ and every normal matrix changes in operator norm by at most
$2M\beta$.  The polynomial conditioning just cited lets us choose
$\beta=(MU/\eps)^{-C}$, for one sufficiently large constant $C$, so that the
induced error is at most $\alpha$.  Maxima and predicates are handled by
Lemma~\ref{lem:hidden-max}; an interval that straddles a threshold may take
either branch, since the two returned coordinate values then differ by at
most the same interface tolerance.  All remaining state maps are compositions
of arithmetic, barrier derivatives, powers, medians, and clipping.  On the
source polynomial envelope their denominators are bounded away from zero by
an inverse polynomial; the smooth maps are therefore polynomially Lipschitz,
while median and clipping are nonexpansive.  Hence their accumulated error
over one batch is $\poly(MU/\eps)\beta$, which is again at most $\alpha$ after
increasing $C$.  The support construction in
Subsection~\ref{subsec:mixed-support} and the leverage construction in the
preceding subsection already give the additive and multiplicative guarantees
required by the proposition.

For a normal solve, freeze $\widehat D$ and its right-hand side as in
\eqref{eq:dyadic-diagonal}, and put
\[
 \widehat L=A^\top\widehat D A,\qquad
 \widehat b=A^\top\widehat Dq.
\]
The sparsifier gives
$\frac12\widehat L\preceq P\preceq2\widehat L$.
Preconditioned Richardson iteration therefore contracts the
$\widehat L$-energy error by a constant factor.  After
$O(\log(1/\alpha))=O(\log(MU/\eps))$ iterations it returns
$\widetilde\pi$ satisfying
\[
 \|\widetilde\pi-\widehat L^{-1}\widehat b\|_{\widehat L}
 \le
 \alpha\|\widehat L^{-1}\widehat b\|_{\widehat L};
\]
when $\widehat b=0$ we return zero.  Multiplication by the full
$\widehat L$, rather than by the sparsifier, makes this a solve of the
desired frozen system.  The perturbation
\eqref{eq:dyadic-diagonal} and polynomial conditioning transfer the same
bound, up to a polynomial factor already absorbed in $C_{\rm stab}$, to the
analytic system in \eqref{eq:exact-newton}.

It remains to explain why the exact repair does not spoil that bound.  Write
$e=\widetilde h-h$, where $h$ is the analytic direction requested by the
stable execution.  Since $A^\top h=0$,
\[
 h_{\rm rep}-h=e-R_\star A^\top e .
\]
In the local norm this gives
\begin{align*}
 \|H^{1/2}(h_{\rm rep}-h)\|_2
 &\le
 \bigl(1+
 \|H^{1/2}R_\star\|_2
 \|A^\top H^{-1/2}\|_2\bigr)
 \|H^{1/2}e\|_2\\
 &\le \poly(MU/\eps)\,\|H^{1/2}e\|_2.
\end{align*}
The first inequality is the triangle inequality followed by
submultiplicativity.  The second follows from the source Hessian bounds and from
$\|A^\top H^{-1/2}\|_2^2=\lambda_{\max}(A^\top H^{-1}A)$.
Thus requesting the raw direction to a correspondingly smaller fixed
inverse-polynomial tolerance makes the repaired direction a valid inexact
direction.  The raw coordinate circuit is formed from the frozen dyadic row
values and the dyadic approximate potential, so
Lemma~\ref{lem:aggregate} computes $A^\top\widetilde h$ exactly with
$r\log^{O(1)}r$ queries.  Therefore
\eqref{eq:star-right-inverse} gives the exact identity
\[
 A^\top h_{\rm rep}
 =A^\top\widetilde h-A^\top R_\star A^\top\widetilde h=0.
\]
Starting from the supplied exactly feasible point, induction keeps every
iterate exactly feasible.  The repair touches only explicit star rows, so it
does not disturb duplicate-copy symmetry or pair-evaluability.

All arithmetic between queries is free in the query model.  The only finite
representations seen by the oracle reductions are the dyadic interfaces, and
their accuracy is inverse polynomial, so their bit length is polylogarithmic
in $MU/\eps$.  Richardson adds only another logarithmic factor.
If a randomized spectral call fails, or an actual row fails its domain check,
the capped execution rejects; on the simultaneous success event neither
occurs.  The imported stability conclusion then gives strict interiority and
the $\eps$ objective guarantee, proving the lemma.
\end{proof}

\paragraph{End of the proof of Theorem~\ref{thm:implicit-lp}.}
Proposition~\ref{prop:ls-batches} uses
$\sqrt r\log^{O(1)}r$ primitive batches.  Lemmas
\ref{lem:transcript-simulation} and~\ref{lem:solver-interface} implement each
batch in $r\log^{O(1)}r$ cut queries, while preserving exact affine
feasibility.  Multiplying the two bounds gives
$r^{3/2}\log^{O(1)}r$.  The conditional success probabilities of the
spectral calls are combined in Section~\ref{sec:completion}.  This proves
Theorem~\ref{thm:implicit-lp}.
\section{Completing the reachability algorithm}
\label{sec:completion}

Apply Theorem~\ref{thm:implicit-lp} to \eqref{eq:phase-one} with output
accuracy
\[
        \eps_{\rm LP}=n^{-30}.
\]
On the successful oracle event, the represented output is in the hidden box,
exactly feasible, and satisfies
\begin{equation}\label{eq:artificial-objective}
        \sum_v(y_v^++y_v^-)\le n^{-30}.
\end{equation}
Exact feasibility gives
\[
 B^\top z-(1-\eta)b=-(y^+-y^-).
\]
Because the artificial variables are nonnegative,
\begin{equation}\label{eq:objective-to-residual}
 \norm{B^\top z-(1-\eta)b}_1
 \le\sum_v|y_v^+-y_v^-|
 \le\sum_v(y_v^++y_v^-)
 \le n^{-30}.
\end{equation}

\paragraph{Acceptance and the final algorithm.}
Run one capped execution of the implicit Phase-I solver.  It returns either a
pair-evaluable candidate $z$ or \textsc{Reject}.  If it rejects, return
\textsc{No}.  Otherwise test \eqref{eq:certificate}.  An out-of-box actual
pair is found by filtered search.  For the residual test, use
the public circuit clipped to $[-1,1]$ (which agrees with $z$ on every actual
pair after the box test) and Lemma~\ref{lem:rounded-aggregate} to enclose every coordinate of
$B^\top z-(1-\eta)b$ within absolute error $n^{-25}$, and accept only if the
resulting upper bound on its $\ell_1$ norm is at most $n^{-20}$.  This uses
$n\log^{O(1)}n$ queries.  Equation~\eqref{eq:objective-to-residual} shows that
a successful solver output is accepted, since its true residual is at most
$n^{-30}$.

For an accepted candidate, construct the threshold digraph $J$ implicitly and
run Lemma~\ref{lem:implicit-bfs} from $s$.  Return \textsc{Yes} exactly when
that search discovers $t$.  Every accepted candidate is deterministically
safe: the conservative test certifies \eqref{eq:certificate}, so
Lemmas~\ref{lem:gap} and~\ref{lem:threshold} show that $J$ has the same
reachability relation as the input.  Returning \textsc{No} is always correct
on a no-instance; on a yes-instance it can be wrong only if the implicit
solver's high-probability event fails.

\paragraph{Query and failure bound.}
Let $N_{\rm spec}$ be a deterministic upper bound on the number of spectral
calls in the single capped execution; by
Proposition~\ref{prop:ls-batches},
$N_{\rm spec}\le\sqrt n\log^{O(1)}n$.  Give each potential call conditional
failure probability at most
$n^{-C_{\rm fail}}/N_{\rm spec}$, conditional on the entire preceding
transcript.  The adaptive union bound is then at most
$n^{-C_{\rm fail}}$.  Conditional on its complement, every primitive meets
the hypotheses of Lemma~\ref{lem:solver-interface}, so the candidate is
returned and accepted.

The LP execution uses $\sqrt n\log^{O(1)}n$ batches of
$n\log^{O(1)}n$ queries.  The acceptance test costs
$n\log^{O(1)}n$, and the final search costs $O(n\log n)$.  Fixed caps on every
spanner attempt, Richardson refinement, and primitive batch give the same
$n^{3/2}\log^{O(1)}n$ upper bound on every random tape.  This proves
Theorem~\ref{thm:main}.

\section{Discussion}

The main open question is whether directed reachability admits
$n^{1+o(1)}$ outgoing-cut queries.  It would also be interesting to obtain
meaningful superlinear lower bounds, to make the pair-evaluable spectral
routine efficient in computation as well as queries, and to remove the
linear dependence on a nonconstant unordered-pair multiplicity cap.

\small\sloppy
\begin{thebibliography}{99}

\bibitem{AnandSaranurakWang2025}
A.~Anand, T.~Saranurak, and Y.~Wang.
\newblock Deterministic edge connectivity and max flow using subquadratic
cut queries.
\newblock In \emph{Proceedings of the 2025 Annual ACM--SIAM Symposium on
Discrete Algorithms (SODA)}, pages 124--142, 2025.
\newblock \url{https://doi.org/10.1137/1.9781611978322.4}.

\bibitem{ApersEtAl2022}
S.~Apers, Y.~Efron, P.~Gawrychowski, T.~Lee, S.~Mukhopadhyay, and
D.~Nanongkai.
\newblock Cut query algorithms with star contraction.
\newblock arXiv:2201.05674, 2022.
\newblock \url{https://arxiv.org/abs/2201.05674}.

\bibitem{BaswanaSen2007}
S.~Baswana and S.~Sen.
\newblock A simple and linear time randomized algorithm for computing sparse
spanners in weighted graphs.
\newblock \emph{Random Structures \& Algorithms}, 30(4):532--563, 2007.
\newblock \url{https://doi.org/10.1002/rsa.20130}.

\bibitem{JiangNanongkaiSawettamalya2026}
Y.~Jiang, D.~Nanongkai, and P.~Sawettamalya.
\newblock Minimum $s$--$t$ cuts with fewer cut queries.
\newblock In \emph{Proceedings of the 2026 Annual ACM--SIAM Symposium on
Discrete Algorithms (SODA)}, pages 258--296, 2026.
\newblock \url{https://doi.org/10.1137/1.9781611978971.12}.

\bibitem{Koutis2014}
I.~Koutis.
\newblock Simple parallel and distributed algorithms for spectral graph
sparsification.
\newblock In \emph{Proceedings of the 26th ACM Symposium on Parallelism in
Algorithms and Architectures (SPAA)}, pages 61--66, 2014.
\newblock \url{https://doi.org/10.1145/2612669.2612676}.

\bibitem{LeeSidfordPathII}
Y.~T.~Lee and A.~Sidford.
\newblock Path Finding II: An $\widetilde O(m\sqrt n)$ algorithm for the
minimum cost flow problem.
\newblock arXiv:1312.6713v2, 2015.
\newblock \url{https://arxiv.org/abs/1312.6713}.

\bibitem{LeeSidford2020}
Y.~T.~Lee and A.~Sidford.
\newblock Solving linear programs with
$\widetilde O(\sqrt{\operatorname{rank}})$ linear system solves.
\newblock arXiv:1910.08033v2, 2020.
\newblock \url{https://arxiv.org/abs/1910.08033}.

\bibitem{NanongkaiOpen}
D.~Nanongkai.
\newblock Cut query for reachability: open problem.
\newblock Open-problem page, first offered at ICALP 2024.
\newblock \url{https://sites.google.com/site/dannanongkai/open}.

\bibitem{RubinsteinSchrammWeinberg2018}
A.~Rubinstein, T.~Schramm, and S.~M.~Weinberg.
\newblock Computing exact minimum cuts without knowing the graph.
\newblock In \emph{9th Innovations in Theoretical Computer Science
Conference (ITCS)}, volume 94 of LIPIcs, pages 39:1--39:16, 2018.
\newblock \url{https://doi.org/10.4230/LIPIcs.ITCS.2018.39}.

\bibitem{Tropp2012}
J.~A.~Tropp.
\newblock User-friendly tail bounds for sums of random matrices.
\newblock \emph{Foundations of Computational Mathematics}, 12(4):389--434,
2012.
\newblock \url{https://doi.org/10.1007/s10208-011-9099-z}.

\end{thebibliography}

\end{document}
